设三场满足命题 \ref{prop:evolution} 的条件, 普通连续谱上 $s_{11}s_{22}\ne0$, 离散谱有限且简单. 记 $s_{22}$ 在 $D_+$ 中的全部零点为 $\zeta_1,\ldots,\zeta_m$, $m=2N$, 范数常数及 $C_j=c_je^{2\ii\theta(\zeta_j)}$ 按第 \ref{ch:discrete} 节定义. 将 $\Sigma$ 定向为 $D_+$ 位于左侧: 实轴的外部两段向右, 内部区间向左, 上半圆顺时针, 下半圆逆时针. 下标 $+$、$-$ 分别表示左、右非切向边界值. 普通轮廓指删去 $0,\pm q_0,\pm\ii q_0,z_a^\pm$ 后的各弧.

\begin{theorem}[非零背景的矩阵 RHP]\label{thm:RHP}\label{prop:jump}
由 Jost 解析列定义
\begin{equation}\label{eq:Mdef}
M(z)=\begin{cases}
M_+(z)=\left(\mu_{-,1},\dfrac{\mu_{+,2}}{s_{22}}\right),&z\in D_+,\\[6pt]
M_-(z)=\left(\dfrac{\mu_{+,1}}{s_{11}},\mu_{-,2}\right),&z\in D_-.
\end{cases}
\end{equation}
则 $M$ 在两域内亚纯, $D_+$ 中仅第二列可在 $\zeta_j$ 处有简单极点, $D_-$ 中仅第一列可在 $\conj{\zeta_j}$ 处有简单极点. 普通轮廓上的跳跃为
\begin{equation}\label{eq:jump}
M_+=M_-J,\qquad
J=\begin{pmatrix}
1&\widetilde\rho\,e^{2\ii\theta}\\
-\rho\,e^{-2\ii\theta}&1-\rho\widetilde\rho
\end{pmatrix},\qquad \det J=1.
\end{equation}
各简单极点的留数满足 \eqref{eq:resplus}--\eqref{eq:resminus}. 行列式、两处归一化与约化为
\begin{align}
\det M(z)&=\gamma(z),\label{eq:detM}\\
M(z)&=I+\dfrac{M^{(1)}}{z}+O(z^{-2}),\qquad z\to\infty,\label{eq:Minfty}\\
M(z)&=-\dfrac{\ii q_0}{z}\sigma_1+O(1),\qquad z\to0,\label{eq:Mzero}\\
M(z)&=\sigma_2\conj{M(\conj{z})}\sigma_2,\qquad
M(h(z))=M(z)\dfrac{\ii z}{q_0}\sigma_1.\label{eq:Msym}
\end{align}
\end{theorem}
\begin{proof}
解析性与列极点由第 \ref{ch:direct}、\ref{ch:scattering}、\ref{ch:discrete} 节的结果得到. 去除散射关系中的指数, 有
\[
\mu_{+,1}=s_{11}\mu_{-,1}+s_{21}e^{-2\ii\theta}\mu_{-,2},\qquad
\mu_{+,2}=s_{12}e^{2\ii\theta}\mu_{-,1}+s_{22}\mu_{-,2}.
\]
第一式给出 $M_{+,1}=M_{-,1}-\rho e^{-2\ii\theta}M_{-,2}$, 代入第二式得到
\[
M_{+,2}=\widetilde\rho e^{2\ii\theta}M_{-,1}+(1-\rho\widetilde\rho)M_{-,2},
\]
即 \eqref{eq:jump}. Wronskian 表达式给出
\[
\det M_+=\dfrac{\det(\mu_{-,1},\mu_{+,2})}{s_{22}}=\gamma,
\]
另一域同理. 第 \ref{ch:asymptotics} 节的展开确定两处归一化. 共轭约化交换两列和两域; 换叶约化将 $\mu_{-,1}$ 与 $\mu_{-,2}$、$\dfrac{\mu_{+,2}}{s_{22}}$ 与 $\dfrac{\mu_{+,1}}{s_{11}}$ 配对, 因而得到 \eqref{eq:Msym}.
\end{proof}

% BEGIN EXPANSION 08-A
\subsubsection*{逐列消元, 跳跃方向与行列式的逐项计算}
这里 $M_{+,1}$ 表示 $M_+$ 的第一列, 不是矩阵的 $(+,1)$ 元.
两个域中的定义逐列写为
\[
\begin{aligned}
M_{+,1}&=\mu_{-,1},&
M_{+,2}&=\dfrac{\mu_{+,2}}{s_{22}},\\
M_{-,1}&=\dfrac{\mu_{+,1}}{s_{11}},&
M_{-,2}&=\mu_{-,2}.
\end{aligned}
\]
第一条散射关系两侧除以非零的 $s_{11}$, 得
\[
\dfrac{\mu_{+,1}}{s_{11}}
=\mu_{-,1}
+\dfrac{s_{21}}{s_{11}}e^{-2\ii\theta}\mu_{-,2}.
\]
把 $\rho=\dfrac{s_{21}}{s_{11}}$ 代入, 再将右侧最后一项移到左边, 就得到
\[
\begin{aligned}
M_{-,1}
&=M_{+,1}+\rho e^{-2\ii\theta}M_{-,2},\\
M_{+,1}
&=M_{-,1}-\rho e^{-2\ii\theta}M_{-,2}.
\end{aligned}
\]
第二条散射关系除以 $s_{22}$, 使用 $\widetilde\rho=\dfrac{s_{12}}{s_{22}}$, 得
\[
M_{+,2}
=\widetilde\rho e^{2\ii\theta}M_{+,1}+M_{-,2}.
\]
将刚得到的第一列式代入, 每一步只对标量与列向量作分配律:
\[
\begin{aligned}
M_{+,2}
&=\widetilde\rho e^{2\ii\theta}
\bigl(M_{-,1}-\rho e^{-2\ii\theta}M_{-,2}\bigr)
+M_{-,2}\\
&=\widetilde\rho e^{2\ii\theta}M_{-,1}
-\widetilde\rho\rho e^{2\ii\theta-2\ii\theta}M_{-,2}
+M_{-,2}\\
&=\widetilde\rho e^{2\ii\theta}M_{-,1}
+(1-\rho\widetilde\rho)M_{-,2}.
\end{aligned}
\]
右乘矩阵的第一列给第一条线性组合, 第二列给第二条线性组合, 因而
\[
\begin{aligned}
M_-J
&=(M_{-,1},M_{-,2})
\begin{pmatrix}
1&\widetilde\rho e^{2\ii\theta}\\
-\rho e^{-2\ii\theta}&1-\rho\widetilde\rho
\end{pmatrix}\\
&=\bigl(
M_{-,1}-\rho e^{-2\ii\theta}M_{-,2},
\widetilde\rho e^{2\ii\theta}M_{-,1}
+(1-\rho\widetilde\rho)M_{-,2}
\bigr)\\
&=(M_{+,1},M_{+,2}).
\end{aligned}
\]
行列式中两项的相消为
\[
\begin{aligned}
\det J
&=1\cdot(1-\rho\widetilde\rho)
-(\widetilde\rho e^{2\ii\theta})
(-\rho e^{-2\ii\theta})\\
&=1-\rho\widetilde\rho+\rho\widetilde\rho=1.
\end{aligned}
\]
由 $\det S=1$ 也直接得到
\[
\begin{aligned}
1
&=s_{11}s_{22}-s_{12}s_{21}\\
&=s_{11}s_{22}
\left(1-\dfrac{s_{12}}{s_{22}}\dfrac{s_{21}}{s_{11}}\right)\\
&=s_{11}s_{22}(1-\rho\widetilde\rho).
\end{aligned}
\]
这解释后面标量迹公式中倒数的来源.
计算 $\det M_+$ 时, 先在第二列使用散射展开:
\[
\begin{aligned}
\det(\mu_{-,1},\mu_{+,2})
&=\det\bigl(\mu_{-,1},
s_{12}e^{2\ii\theta}\mu_{-,1}+s_{22}\mu_{-,2}\bigr)\\
&=s_{12}e^{2\ii\theta}
\det(\mu_{-,1},\mu_{-,1})
+s_{22}\det(\mu_{-,1},\mu_{-,2})\\
&=0+s_{22}\gamma.
\end{aligned}
\]
同样, 对 $M_-$ 的第一列展开:
\[
\begin{aligned}
\det(\mu_{+,1},\mu_{-,2})
&=s_{11}\det(\mu_{-,1},\mu_{-,2})
+s_{21}e^{-2\ii\theta}\det(\mu_{-,2},\mu_{-,2})\\
&=s_{11}\gamma.
\end{aligned}
\]
两侧分别除以相应的透射系数, 都得到 $\gamma$.
在离散零点处, 该商恒等式由亚纯延拓成立, 不能把零点处的表达式理解为直接作 $\dfrac{0}{0}$ 数值除法.

\paragraph{标准基与秩一主部.}
标准基及其转置为
\[
e_1=\begin{pmatrix}1\\0\end{pmatrix},\qquad
e_2=\begin{pmatrix}0\\1\end{pmatrix},\qquad
e_1^T=(1,0),\qquad e_2^T=(0,1).
\]
因而两种留数矩阵按元展开是
\[
\begin{aligned}
C_j\mathbf v_je_2^T
&=\begin{pmatrix}
0&C_j(\mathbf v_j)_1\\
0&C_j(\mathbf v_j)_2
\end{pmatrix},\\
-\conj{C_j}\mathbf u_je_1^T
&=\begin{pmatrix}
-\conj{C_j}(\mathbf u_j)_1&0\\
-\conj{C_j}(\mathbf u_j)_2&0
\end{pmatrix}.
\end{aligned}
\]
第一种主部只占第二列, 所以在 $z=\zeta_j$ 取第一列值没有奇性.
第二种主部只占第一列, 所以在 $z=\conj{\zeta_j}$ 取第二列值没有奇性.
这正是后面积分--代数系统可以在极点位置求值的原因.
% END EXPANSION 08-A
\subsection{留数条件与重构公式}
记 $e_1,e_2$ 为标准基, 置
\begin{equation}\label{eq:Gdefinition}
G=I-J^{-1}=\begin{pmatrix}
\rho\widetilde\rho&\widetilde\rho e^{2\ii\theta}\\
-\rho e^{-2\ii\theta}&0
\end{pmatrix},\qquad M_+-M_-=M_+G.
\end{equation}
以下表示按照先消去奇性、再应用 Plemelj 公式的顺序建立\cite{ZhangFan2020,Fan2022}. 它适用于满足所列条件的 RHP 解.

% BEGIN EXPANSION 08-B
\paragraph{先核对逆矩阵, 再核对 $G$ 的符号.}
因为 $\det J=1$, 二阶逆矩阵公式给出
\[
J^{-1}=\begin{pmatrix}
1-\rho\widetilde\rho&-\widetilde\rho e^{2\ii\theta}\\
\rho e^{-2\ii\theta}&1
\end{pmatrix}.
\]
其乘积四个元逐个为
\[
\begin{aligned}
(JJ^{-1})_{11}
&=1-\rho\widetilde\rho
+\widetilde\rho e^{2\ii\theta}\rho e^{-2\ii\theta}=1,\\
(JJ^{-1})_{12}
&=-\widetilde\rho e^{2\ii\theta}
+\widetilde\rho e^{2\ii\theta}=0,\\
(JJ^{-1})_{21}
&=-\rho e^{-2\ii\theta}(1-\rho\widetilde\rho)
+(1-\rho\widetilde\rho)\rho e^{-2\ii\theta}=0,\\
(JJ^{-1})_{22}
&=\rho\widetilde\rho+1-\rho\widetilde\rho=1.
\end{aligned}
\]
再按元作 $I-J^{-1}$, 得
\[
\begin{aligned}
(I-J^{-1})_{11}&=1-(1-\rho\widetilde\rho)=\rho\widetilde\rho,\\
(I-J^{-1})_{12}&=0-(-\widetilde\rho e^{2\ii\theta})
=\widetilde\rho e^{2\ii\theta},\\
(I-J^{-1})_{21}&=0-\rho e^{-2\ii\theta}
=-\rho e^{-2\ii\theta},\\
(I-J^{-1})_{22}&=1-1=0.
\end{aligned}
\]
从 $M_+=M_-J$ 右乘 $J^{-1}$ 得 $M_-=M_+J^{-1}$, 所以
\[
M_+-M_-=M_+-M_+J^{-1}=M_+(I-J^{-1})=M_+G.
\]
这里右侧是 $M_+$, 而不是 $M_-$.
若将密度错误地写成 $M_-(I-J^{-1})$, 就不能重现原跳跃.
按列计算实际密度, 有
\[
\begin{aligned}
W_1
&=\rho\widetilde\rho M_{+,1}
-\rho e^{-2\ii\theta}M_{+,2},\\
W_2
&=\widetilde\rho e^{2\ii\theta}M_{+,1}.
\end{aligned}
\]
尤其
\[
W_{12}=\widetilde\rho e^{2\ii\theta}(M_{+,1})_1.
\]
这些等式只展开原有密度, 没有把它当作预先已知的独立函数.
因此含反射表示仍是未知 $M_+$ 与离散正则列共同满足的关系.
% END EXPANSION 08-B
\begin{proposition}[含反射的留数积分表示]\label{prop:cauchy-representation}
设定理 \ref{thm:RHP} 的 RHP 有解. 假设减去离散主部及 $z=0$ 主部后, 矩阵在各有限特殊点有有界、相容的扇区边界值; 跳跃密度 $W=M_+G$ 在普通弧上局部 Hölder 连续, 且 $W,sW\in L^1(\Sigma)$. 删去特殊点的小圆弧时, 假设其上的正则化矩阵一致有界. 则
\begin{equation}\label{eq:general-cauchy}
\begin{aligned}
M(z)={}&I-\dfrac{\ii q_0}{z}\sigma_1
+\sum\limits_{j=1}^{m}\left(
\dfrac{C_j\mathbf v_je_2^T}{z-\zeta_j}
-\dfrac{\conj{C_j}\mathbf u_je_1^T}{z-\conj{\zeta_j}}\right)\\
&+\dfrac{1}{2\pi\ii}\int\limits_\Sigma\dfrac{M_+(s)G(s)}{s-z}\,\dd s,
\qquad z\in\mathbb C\mathbin{\backslash}\Sigma,
\end{aligned}
\end{equation}
其中 $\mathbf v_j=M_1(\zeta_j)$、$\mathbf u_j=M_2(\conj{\zeta_j})$ 是各点的正则列. 它们满足积分--代数关系
\begin{align}
\mathbf v_j={}&e_1-\dfrac{\ii q_0}{\zeta_j}e_2
-\sum\limits_{k=1}^{m}\dfrac{\conj{C_k}\mathbf u_k}{\zeta_j-\conj{\zeta_k}}
+\dfrac{1}{2\pi\ii}\int\limits_\Sigma\dfrac{W_1(s)}{s-\zeta_j}\,\dd s,
\label{eq:general-finite1}\\
\mathbf u_j={}&e_2-\dfrac{\ii q_0}{\conj{\zeta_j}}e_1
+\sum\limits_{k=1}^{m}\dfrac{C_k\mathbf v_k}{\conj{\zeta_j}-\zeta_k}
+\dfrac{1}{2\pi\ii}\int\limits_\Sigma\dfrac{W_2(s)}{s-\conj{\zeta_j}}\,\dd s.
\label{eq:general-finite2}
\end{align}
按 $q=\ii M^{(1)}_{12}$ 定义重构光场, 则
\begin{equation}\label{eq:qgeneral}
q=q_0+\ii\sum\limits_{j=1}^{m}C_j(\mathbf v_j)_1
-\dfrac{1}{2\pi}\int\limits_\Sigma W_{12}(s)\,\dd s.
\end{equation}
\end{proposition}
\begin{proof}
第 \ref{ch:discrete} 节的留数条件给出
\[
\Res_{z=\zeta_j}M=C_j\mathbf v_je_2^T,\qquad
\Res_{z=\conj{\zeta_j}}M=-\conj{C_j}\mathbf u_je_1^T.
\]
从 $M$ 减去 $I-\dfrac{\ii q_0}{z}\sigma_1$ 及上述各主部, 所得分片解析矩阵 $F_0$ 在无穷远为 $O(z^{-1})$, 而 $F_{0,+}-F_{0,-}=W$. 对删去特殊点小圆盘并截断于大圆的两域分别使用 Cauchy 公式. 左侧为正向边界, 将两域的公共边界合并得到 $W$; 局部一致有界性保证小圆贡献消失, 大圆上的 $O(z^{-1})$ 使其贡献消失. 因而
\[
F_0(z)=\dfrac{1}{2\pi\ii}\int\limits_\Sigma\dfrac{W(s)}{s-z}\,\dd s,
\]
得到 \eqref{eq:general-cauchy}. 正号由 \eqref{eq:Gdefinition} 与轮廓定向共同确定.

在第一列取 $z=\zeta_j$, 在第二列取 $z=\conj{\zeta_j}$, 各自本列在该点没有离散极点, 得 \eqref{eq:general-finite1}--\eqref{eq:general-finite2}. 沿避开轮廓的闭角域令 $z\to\infty$, 使用
\[
\dfrac{1}{s-z}=-\dfrac{1}{z}+\dfrac{s}{z(s-z)},\qquad
\left|\dfrac{s}{z(s-z)}\right|\le\dfrac{C|s|}{|z|^2},
\]
及 $sW\in L^1$ 控制余项, 得
\begin{equation}\label{eq:M1-general}
M^{(1)}=-\ii q_0\sigma_1
+\sum\limits_{j=1}^{m}(C_j\mathbf v_je_2^T-\conj{C_j}\mathbf u_je_1^T)
-\dfrac{1}{2\pi\ii}\int\limits_\Sigma W(s)\,\dd s.
\end{equation}
读取 $(1,2)$ 元的系数并乘以 $\ii$, 即得 \eqref{eq:qgeneral}. 对直接散射构造的 $M$, 此定义与原光场一致, 见下述重构命题.
\end{proof}

% BEGIN EXPANSION 08-C
\paragraph{Cauchy 表示中每一种奇性的消去.}
原证明使用的正则化矩阵完整写为
\[
\begin{aligned}
F_0(z)
={}&M(z)-I+\dfrac{\ii q_0}{z}\sigma_1\\
&-\sum\limits_{j=1}^{m}
\dfrac{C_j\mathbf v_je_2^T}{z-\zeta_j}
+\sum\limits_{j=1}^{m}
\dfrac{\conj{C_j}\mathbf u_je_1^T}{z-\conj{\zeta_j}}.
\end{aligned}
\]
在 $\zeta_j$ 附近, $M$ 的 Laurent 展开首项是
$\dfrac{C_j\mathbf v_je_2^T}{z-\zeta_j}$.
上式中减去完全相同的项, 所以 $F_0$ 没有该极点.
在 $\conj{\zeta_j}$ 附近首项是
$-\dfrac{\conj{C_j}\mathbf u_je_1^T}{z-\conj{\zeta_j}}$;
这里在 $F_0$ 中加上其相反数, 同样消去极点.
在 $0$ 附近, 原归一化是
\[
M(z)=-\dfrac{\ii q_0}{z}\sigma_1+O(1),
\]
所以 $M+\dfrac{\ii q_0}{z}\sigma_1$ 有界.
在无穷远, 每个被减去的离散主部都是 $O(z^{-1})$, 故
\[
F_0(z)=O(z^{-1}).
\]
所有被减去的有理项在普通弧两侧取值一致, 所以
\[
\begin{aligned}
F_{0,+}-F_{0,-}
&=(M_+-M_-)+0+0+0\\
&=W.
\end{aligned}
\]
对任一不在轮廓上的固定 $z$, 特殊点周围半径为 $\varepsilon$ 的小圆弧上,
$\dfrac{1}{|s-z|}$ 有统一上界.
若 $F_0$ 在这些弧上一致有界, 则弧长为 $O(\varepsilon)$ 给出
\[
\left\|
\int_{\text{小圆弧}}
\dfrac{F_0(s)}{s-z}\,\dd s
\right\|
\le
\sup\limits_{\text{小圆弧}}\|F_0(s)\|
\sup\limits_{\text{小圆弧}}\dfrac{1}{|s-z|}
\operatorname{length}(\text{小圆弧})
\longrightarrow0.
\]
在半径为 $R$ 的大圆上, $F_0(s)=O(R^{-1})$ 且
$\dfrac{1}{s-z}=O(R^{-1})$, 从而
\[
\int_{|s|=R}\dfrac{F_0(s)}{s-z}\,\dd s=O(R^{-1})\longrightarrow0.
\]
沿公共边界, $D_+$ 的正向边界与 $\Sigma$ 同向,
$D_-$ 的正向边界与 $\Sigma$ 反向.
两条边界积分相加, 密度因此是 $F_{0,+}-F_{0,-}=W$, 不是其相反数.

\paragraph{边界值与积分方程的核对.}
在普通弧的局部光滑点, Plemelj 公式按既定左侧正向给出
\[
\begin{aligned}
F_{0,+}(z)
&=\dfrac{1}{2}W(z)
+\dfrac{1}{2\pi\ii}\operatorname{p.v.}
\int\limits_\Sigma\dfrac{W(s)}{s-z}\,\dd s,\\
F_{0,-}(z)
&=-\dfrac{1}{2}W(z)
+\dfrac{1}{2\pi\ii}\operatorname{p.v.}
\int\limits_\Sigma\dfrac{W(s)}{s-z}\,\dd s.
\end{aligned}
\]
相减恰得 $W$.
把第一式加回全部有理项, 就给出 $M_+$ 的边界积分方程;
再用 $W=M_+G$, 未知边界值同时出现在积分外和积分内.
这一推导说明原命题是已有 RHP 解的必要表示.
要从该式独立证明一般含反射 RHP 有解, 还需要解相应的奇异积分方程;
这里没有加入这一额外存在性结论.

\paragraph{在正则列取值时, 为什么没有自作用发散项.}
因为
\[
e_2^Te_1=0,\quad e_1^Te_1=1,\quad
e_1^Te_2=0,\quad e_2^Te_2=1,
\]
对表示式右乘 $e_1$, 得
\[
M_1(z)=e_1-\dfrac{\ii q_0}{z}e_2
-\sum\limits_{k=1}^{m}
\dfrac{\conj{C_k}\mathbf u_k}{z-\conj{\zeta_k}}
+\dfrac{1}{2\pi\ii}\int\limits_\Sigma
\dfrac{W_1(s)}{s-z}\,\dd s.
\]
所有 $D_+$ 极点主部都因 $e_2^Te_1=0$ 消失.
于是直接令 $z=\zeta_j$ 就是 \eqref{eq:general-finite1}.
右乘 $e_2$ 则得到
\[
M_2(z)=e_2-\dfrac{\ii q_0}{z}e_1
+\sum\limits_{k=1}^{m}
\dfrac{C_k\mathbf v_k}{z-\zeta_k}
+\dfrac{1}{2\pi\ii}\int\limits_\Sigma
\dfrac{W_2(s)}{s-z}\,\dd s.
\]
所有 $D_-$ 极点主部因 $e_1^Te_2=0$ 消失.
于是令 $z=\conj{\zeta_j}$ 就是 \eqref{eq:general-finite2}.
这里不同域的离散点互异, 所以所剩分母均不为零.

\paragraph{无穷远展开与光场前面的积分负号.}
精确恒等式
\[
\dfrac{1}{s-z}+\dfrac{1}{z}
=\dfrac{z+s-z}{z(s-z)}
=\dfrac{s}{z(s-z)}
\]
给出
\[
\begin{aligned}
\dfrac{1}{2\pi\ii}\int\limits_\Sigma
\dfrac{W(s)}{s-z}\,\dd s
={}&-\dfrac{1}{2\pi\ii z}\int\limits_\Sigma W(s)\,\dd s\\
&+\dfrac{1}{2\pi\ii}\int\limits_\Sigma
\dfrac{sW(s)}{z(s-z)}\,\dd s.
\end{aligned}
\]
后项范数不超过常数乘
\[
\dfrac{1}{|z|^2}\int\limits_\Sigma |s|\|W(s)\|\,|\dd s|,
\]
所以是 $O(z^{-2})$.
离散有理项逐项使用
\[
\dfrac{1}{z-\zeta_j}=\dfrac{1}{z}+O(z^{-2}),\qquad
\dfrac{1}{z-\conj{\zeta_j}}=\dfrac{1}{z}+O(z^{-2}).
\]
因此 $(1,2)$ 元的系数为
\[
M^{(1)}_{12}
=-\ii q_0
+\sum\limits_{j=1}^{m}C_j(\mathbf v_j)_1
-\dfrac{1}{2\pi\ii}\int\limits_\Sigma W_{12}(s)\,\dd s.
\]
第二组留数矩阵的第二列是零, 对这个矩阵元不贡献.
两侧乘以 $\ii$, 用 $\ii(-\ii)=1$ 及
$\dfrac{\ii}{2\pi\ii}=\dfrac{1}{2\pi}$, 就得到 \eqref{eq:qgeneral}.
% END EXPANSION 08-C
\begin{proposition}[光场与介质重构]\label{prop:reconstruction}
直接散射构造的矩阵 $M$ 满足
\begin{equation}\label{eq:reconq}
Q(q)=\dfrac{\ii}{2}[\sg,M^{(1)}],\qquad q=\ii M^{(1)}_{12}.
\end{equation}
若满足命题 \ref{prop:resonance} 的共振相容条件, 则 $M$ 在两个共振原像的空间归一化值可逆, 且
\begin{equation}\label{eq:RfromM}
R=M(x,t,z_a)\dfrac{\eta_0\lambda(z_a)}{a}\sg M(x,t,z_a)^{-1},
\qquad z_a=z_a^\pm.
\end{equation}
两式给出同一个介质矩阵. 原物理场由
\begin{equation}\label{eq:original-fields}
\widetilde q=e^{-\ii\Omega t}q,\qquad
\widetilde p=-e^{-\ii\Omega t}R_{12},\qquad \eta=R_{11}
\end{equation}
恢复.
\end{proposition}
\begin{proof}
由 \eqref{eq:mu-highenergy} 与 \eqref{eq:transmission-highenergy}, 在 $D_+$ 中第二列除以 $s_{22}$ 后, 其对角系数从 $-\ii J_+$ 变为 $-\ii(J_++J_{\rm tot})=-\ii J_-$. 另一域相同, 故
\[
M^{(1)}=-\ii B(q)+\ii J_-\sg,
\]
得到 \eqref{eq:reconq}. 这一关系也可由空间方程
\begin{equation}\label{eq:Mspace}
M_x=-\ii k\sg M+QM+\ii\lambda M\sg
\end{equation}
比较无穷远的常数项得到: 第 \ref{ch:asymptotics} 节的积分估计还给出 $M_x=O(z^{-1})$, 而 $k,\lambda\sim\dfrac{z}{2}$.

在共振点, $S_{\rm sp}$ 对角, 两种拼接均为 $M=\mu_-$. 命题 \ref{prop:resonance} 的证明给出
\[
\chi_-^{-1}R\chi_-=\dfrac{\eta_0\lambda(z_a)}{a}\sg.
\]
空间指数与 $\sg$ 对易, 因而得到 \eqref{eq:RfromM}; 这里使用空间归一化值, 不对奇异的传播指数求值. $\gamma(z_a)>0$ 保证可逆. 换叶时 $\lambda$ 变号, 同时 $\sigma_1\sg\sigma_1=-\sg$, 两次变号相消, 故两原像的结果一致.
\end{proof}

% BEGIN EXPANSION 08-D
\subsubsection*{从高能系数到三个场的完整矩阵运算}
原有记号 $B(q)$ 的非对角元是 $q,\conj q$, 因而
\[
M^{(1)}=
\begin{pmatrix}
\ii J_-&-\ii q\\
-\ii\conj q&-\ii J_-
\end{pmatrix}.
\]
分别左乘和右乘 $\sg$:
\[
\sg M^{(1)}=
\begin{pmatrix}
\ii J_-&-\ii q\\
\ii\conj q&\ii J_-
\end{pmatrix},\qquad
M^{(1)}\sg=
\begin{pmatrix}
\ii J_-&\ii q\\
-\ii\conj q&\ii J_-
\end{pmatrix}.
\]
相减并乘 $\dfrac{\ii}{2}$, 得
\[
\dfrac{\ii}{2}[\sg,M^{(1)}]
=\dfrac{\ii}{2}
\begin{pmatrix}
0&-2\ii q\\
2\ii\conj q&0
\end{pmatrix}
=\begin{pmatrix}
0&q\\-\conj q&0
\end{pmatrix}=Q(q).
\]
直接从空间方程也能读出完全相同的式子.
在无穷远有
\[
k=\dfrac{z}{2}-\dfrac{q_0^2}{2z},\qquad
\lambda=\dfrac{z}{2}+\dfrac{q_0^2}{2z},\qquad
M=I+\dfrac{M^{(1)}}{z}+O(z^{-2}).
\]
将右侧逐项展开, 常数项为
\[
\begin{aligned}
-\ii k\sg M
&=-\dfrac{\ii z}{2}\sg
-\dfrac{\ii}{2}\sg M^{(1)}+O(z^{-1}),\\
QM&=Q+O(z^{-1}),\\
\ii\lambda M\sg
&=\dfrac{\ii z}{2}\sg
+\dfrac{\ii}{2}M^{(1)}\sg+O(z^{-1}).
\end{aligned}
\]
因为左侧 $M_x=O(z^{-1})$, 常数项必须为零, 即
\[
0=Q-\dfrac{\ii}{2}\sg M^{(1)}
+\dfrac{\ii}{2}M^{(1)}\sg.
\]
移项便是 \eqref{eq:reconq}.

\paragraph{共振重构的逆矩阵与换叶抵消.}
在共振原像处写出 $2\times2$ 逆矩阵而不省略分母:
\[
M(z_a)^{-1}
=\dfrac{1}{\gamma(z_a)}
\begin{pmatrix}
M_{22}(z_a)&-M_{12}(z_a)\\
-M_{21}(z_a)&M_{11}(z_a)
\end{pmatrix}.
\]
因此重构矩阵四个元可以直接检查为
\[
\begin{aligned}
R_{11}
&=\dfrac{\eta_0\lambda(z_a)}{a\gamma(z_a)}
\bigl(M_{11}(z_a)M_{22}(z_a)+M_{12}(z_a)M_{21}(z_a)\bigr),\\
R_{12}
&=-\dfrac{2\eta_0\lambda(z_a)}{a\gamma(z_a)}
M_{11}(z_a)M_{12}(z_a),\\
R_{21}
&=\dfrac{2\eta_0\lambda(z_a)}{a\gamma(z_a)}
M_{21}(z_a)M_{22}(z_a),\\
R_{22}&=-R_{11}.
\end{aligned}
\]
这里没有在 $\phi$ 的奇异时间指数中代入共振点;
所用对象始终是空间归一化的 $M(z_a)$.
由于 $z_a$ 是非零实数,
\[
\gamma(z_a)=1+\dfrac{q_0^2}{z_a^2}>0.
\]
换叶公式中的标量 $\dfrac{\ii z_a}{q_0}$ 在共轭乘积两侧抵消, 于是
\[
\begin{aligned}
M(h(z_a))\lambda(h(z_a))\sg M(h(z_a))^{-1}
&=M(z_a)\sigma_1(-\lambda(z_a)\sg)\sigma_1M(z_a)^{-1}\\
&=M(z_a)\lambda(z_a)\sg M(z_a)^{-1}.
\end{aligned}
\]
所用反对易关系按元可见:
\[
\sigma_1\sg\sigma_1
=\begin{pmatrix}0&1\\1&0\end{pmatrix}
\begin{pmatrix}1&0\\0&-1\end{pmatrix}
\begin{pmatrix}0&1\\1&0\end{pmatrix}
=\begin{pmatrix}-1&0\\0&1\end{pmatrix}=-\sg.
\]
原物理场重构中的负号来自 $R_{12}=-p$, 所以
$\widetilde p=e^{-\ii\Omega t}p=-e^{-\ii\Omega t}R_{12}$.

% END EXPANSION 08-D
\subsection{迹公式与相位条件}
相位条件由 $z=0$ 处的归一化确定. 透射系数同时包含离散零点和连续谱反射的贡献. 以下标量分解采用非零边界逆散射中的 Cauchy 积分方法\cite{Fan2022,ZhangFan2020}, 积分轮廓取 HMB 的实轴与圆周.

\begin{proposition}[迹公式与同相约束]\label{prop:theta}
设透射系数在普通轮廓上非零, 只有上述有限简单离散零点. 取
$\ell(s)=\log(1-\rho(s)\widetilde\rho(s))$ 的零指标分支, 设其在普通弧上局部 Hölder 连续, 在 $0,\infty$ 的极限为零, 且 $\ell,\dfrac{\ell(s)}{s}\in L^1(\Sigma)$. 假设剔除零点后两域中的透射函数具有与乘积关系相容的单值对数, 这些对数与下述 Cauchy 积分在各有限特殊点的边界值有界并局部匹配. 则
\begin{equation}\label{eq:trace}
s_{22}(z)=\prod_{j=1}^{m}\dfrac{z-\zeta_j}{z-\conj{\zeta_j}}
\exp\left\{-\dfrac{1}{2\pi\ii}\int\limits_\Sigma\dfrac{\ell(s)}{s-z}\,\dd s\right\},\qquad z\in D_+.
\end{equation}
$s_{11}$ 由共轭约化确定. 相同左右背景要求
\begin{equation}\label{eq:theta-condition}
\prod_{j=1}^{m}\dfrac{\zeta_j}{\conj{\zeta_j}}
\exp\left\{-\dfrac{1}{2\pi\ii}\int\limits_\Sigma\dfrac{\ell(s)}{s}\,\dd s\right\}=1.
\end{equation}
特别地, 无反射时 $s_{22}$ 等于 \eqref{eq:trace} 中的有理乘积, 且 $\displaystyle\prod_{j=1}^{m}\dfrac{\zeta_j}{\conj{\zeta_j}}=1$.
\end{proposition}
\begin{proof}
记 $B_0(z)=\displaystyle\prod_{j=1}^{m}\dfrac{z-\zeta_j}{z-\conj{\zeta_j}}$, 定义
$\beta_+=\dfrac{s_{22}}{B_0}$、$\beta_-=s_{11}B_0$. 它们在相应域解析且无零点, 无穷远归一化为 1. 由 $\det S=1$,
\[
\beta_+\beta_-=s_{11}s_{22}=\dfrac{1}{1-\rho\widetilde\rho}.
\]
相容的对数分支满足 $\log\beta_++\log\beta_-=-\ell$, 因而分片解析函数 $f_+=\log\beta_+$、$f_-=-\log\beta_-$ 的跳跃为 $f_+-f_-=-\ell$. Plemelj 公式及归一化给出
\[
f(z)=-\dfrac{1}{2\pi\ii}\int\limits_\Sigma\dfrac{\ell(s)}{s-z}\,\dd s.
\]
差函数在特殊点可去且在无穷远为零, 因而没有附加整函数项. 代回 $\beta_+$ 即得 \eqref{eq:trace}. 为取零点极限, 沿 $D_+$ 内的 $z=-\ii\varepsilon$, $\varepsilon\downarrow0$. 实轴上 $\left|\dfrac{s}{s+\ii\varepsilon}\right|\le1$, 圆周与零点分离, 故 $\dfrac{\ell(s)}{s-z}$ 受 $C\left|\dfrac{\ell(s)}{s}\right|$ 控制. 支配收敛结合 $s_{22}(0)=1$ 给出 \eqref{eq:theta-condition}.
\end{proof}

% BEGIN EXPANSION 08-K
\paragraph{去零点乘积与标量跳跃的符号.}
在 $D_+$ 中, $B_0$ 与 $s_{22}$ 在同一点各有一阶零;
在 $\zeta_j$ 附近,
\[
\begin{aligned}
s_{22}(z)&=s_{22}'(\zeta_j)(z-\zeta_j)+O((z-\zeta_j)^2),\\
B_0(z)&=(z-\zeta_j)
\dfrac{1}{z-\conj{\zeta_j}}
\prod\limits_{\substack{k=1\\k\ne j}}^{m}
\dfrac{z-\zeta_k}{z-\conj{\zeta_k}}.
\end{aligned}
\]
所以其商的有限非零极限为
\[
\beta_+(\zeta_j)
=s_{22}'(\zeta_j)(\zeta_j-\conj{\zeta_j})
\prod\limits_{\substack{k=1\\k\ne j}}^{m}
\dfrac{\zeta_j-\conj{\zeta_k}}{\zeta_j-\zeta_k}.
\]
同理, $D_-$ 中 $s_{11}$ 的零与 $B_0$ 的极点相消,
因此 $\beta_-=s_{11}B_0$ 也解析且无零点.
相乘时 $B_0$ 完全约掉:
\[
\beta_+\beta_-
=\dfrac{s_{22}}{B_0}s_{11}B_0
=s_{11}s_{22}
=\dfrac{1}{1-\rho\widetilde\rho}.
\]
在原文明确假设的相容单值对数分支下, 才能写
\[
\log\beta_++\log\beta_-=-\ell.
\]
该分支条件排除了在各连通域之间任意加上不同 $2\pi\ii$ 整数倍的歧义.
原文将下侧函数定义为 $f_-=-\log\beta_-$, 所以
\[
f_+-f_-=log\beta_++\log\beta_-=-\ell.
\]
故 Cauchy 核前面是负号.
把 $f$ 减去其 Cauchy 积分后, 普通轮廓两侧没有跳跃;
特殊点有界且相容, 所以这些点可去.
无穷远极限为零, 由 Liouville 定理差函数恒为零.
最后在 $D_+$ 中逐步指数化:
\[
\begin{aligned}
\log\beta_+(z)
&=-\dfrac{1}{2\pi\ii}\int\limits_\Sigma
\dfrac{\ell(s)}{s-z}\,\dd s,\\
\beta_+(z)
&=\exp\left\{-\dfrac{1}{2\pi\ii}\int\limits_\Sigma
\dfrac{\ell(s)}{s-z}\,\dd s\right\},\\
s_{22}(z)&=B_0(z)\beta_+(z).
\end{aligned}
\]

\paragraph{零点归一化如何成为同相条件.}
取 $z=-\ii\varepsilon$ 且 $0<\varepsilon<q_0$.
这时 $\operatorname{Im}z<0$ 且 $1-\dfrac{q_0^2}{|z|^2}<0$,
故 $z\in D_+$.
对实轴上的 $s$,
\[
\left|\dfrac{s}{s+\ii\varepsilon}\right|
=\dfrac{|s|}{\sqrt{s^2+\varepsilon^2}}\le1,
\]
从而
\[
\left|\dfrac{\ell(s)}{s+\ii\varepsilon}\right|
\le\left|\dfrac{\ell(s)}{s}\right|.
\]
圆周满足 $|s|=q_0$, 取 $\varepsilon<\dfrac{q_0}{2}$ 时
$|s+\ii\varepsilon|\ge\dfrac{q_0}{2}$,
因此圆周上的同类控制常数可取为 2.
原假设使该控制函数可积, 支配收敛给出
\[
\lim\limits_{\varepsilon\downarrow0}
\int\limits_\Sigma\dfrac{\ell(s)}{s+\ii\varepsilon}\,\dd s
=\int\limits_\Sigma\dfrac{\ell(s)}{s}\,\dd s.
\]
离散乘积逐因子取极限:
\[
\dfrac{-\ii\varepsilon-\zeta_j}
{-\ii\varepsilon-\conj{\zeta_j}}
\longrightarrow\dfrac{\zeta_j}{\conj{\zeta_j}}.
\]
用 $s_{22}(0)=1$ 即得到 \eqref{eq:theta-condition}.
一个四点轨道的另一个 $D_+$ 点是
$-\dfrac{q_0^2}{\conj{\zeta_n}}$, 它的相位比值为
\[
\dfrac{-\dfrac{q_0^2}{\conj{\zeta_n}}}{-\dfrac{q_0^2}{\zeta_n}}
=\dfrac{\zeta_n}{\conj{\zeta_n}}.
\]
两个因子相乘便是 $(\dfrac{\zeta_n}{\conj{\zeta_n}})^2$.
又 $\dfrac{\zeta_n}{\conj{\zeta_n}}=e^{2\ii\arg\zeta_n}$,
所以所有轨道乘积为
\[
\exp\left\{4\ii\sum\limits_{n=1}^{N}\arg\zeta_n\right\}.
\]
等于 1 恰是原文的模 $2\pi$ 相位条件.
% END EXPANSION 08-K
每个完整四点轨道取外部上半平面的代表 $\zeta_n$, 则它的两个 $D_+$ 零点对无反射相位乘积的贡献为 $\left(\dfrac{\zeta_n}{\conj{\zeta_n}}\right)^2$. 因而同相约束化为
\[
4\sum\limits_{n=1}^{N}\arg\zeta_n=0\pmod{2\pi}.
\]
纯虚代表逐轨道满足这一条件.

\subsection{HMB 方程的 \texorpdfstring{$N$}{N} 孤子解}
这里的 $N$ 按完整简单四点轨道的数目计数. 以下在无反射条件下求出矩阵解及相应场的重构公式.

当 $\rho=\widetilde\rho=0$ 时, 跳跃矩阵为 $I$. 在非极点轮廓处取可去条件. 在分支点 $z_b=\pm\ii q_0$ 处, 矩阵 $M$ 正则, 而其行列式 $\gamma(z)=1+\dfrac{q_0^2}{z^2}$ 具有一阶零点, 即 $\gamma(z_b)=0$、$\gamma'(z_b)\ne0$.

固定 $N$ 个两两不交的完整简单四点轨道. 按 \eqref{eq:discrete-indexing} 选取外部上半平面代表 $\zeta_1,\ldots,\zeta_N$, 并置
\begin{equation}\label{eq:N-orbit-data}
\begin{aligned}
m&=2N,\qquad \zeta_{N+n}=-\dfrac{q_0^2}{\conj{\zeta_n}},\qquad n=1,\ldots,N,\\
c_{N+n}&=-\dfrac{q_0^2}{(\conj{\zeta_n})^2}\conj{c_n},\qquad
C_j=c_je^{2\ii\theta(\zeta_j)},\qquad j=1,\ldots,m.
\end{aligned}
\end{equation}
所有 $4N$ 个谱点互异且位于 $\Sigma$ 外, 各 $c_n\ne0$. 同相背景要求
\begin{equation}\label{eq:N-orbit-phase}
\prod\limits_{j=1}^{m}\dfrac{\zeta_j}{\conj{\zeta_j}}=1.
\end{equation}
下文保留全部 $m$ 个 $D_+$ 点及其共轭点, 令 $\mathbf v_j=M_1(\zeta_j)$、$\mathbf u_j=M_2(\conj{\zeta_j})$.

\begin{proposition}[留数系统]\label{prop:finite}
除 $0$ 与指定离散极点外没有有限奇点的无反射矩阵解具有表示
\begin{equation}\label{eq:partialfraction}
M(z)=I-\dfrac{\ii q_0}{z}\sigma_1+\sum\limits_{j=1}^{m}
\left(\dfrac{C_j\mathbf v_je_2^T}{z-\zeta_j}
-\dfrac{\conj{C_j}\mathbf u_je_1^T}{z-\conj{\zeta_j}}\right),
\end{equation}
其正则列满足
\begin{align}
\mathbf v_j&=e_1-\dfrac{\ii q_0}{\zeta_j}e_2
-\sum\limits_{k=1}^{m}\dfrac{\conj{C_k}\mathbf u_k}{\zeta_j-\conj{\zeta_k}},\label{eq:finite1}\\
\mathbf u_j&=e_2-\dfrac{\ii q_0}{\conj{\zeta_j}}e_1
+\sum\limits_{k=1}^{m}\dfrac{C_k\mathbf v_k}{\conj{\zeta_j}-\zeta_k}.\label{eq:finite2}
\end{align}
反之, 任一系统解通过 \eqref{eq:partialfraction} 给出具有指定留数和两处归一化的矩阵. 记
\[
A_{jk}=\dfrac{\conj{C_k}}{\zeta_j-\conj{\zeta_k}},\qquad
B_{jk}=\dfrac{C_k}{\conj{\zeta_j}-\zeta_k},\qquad j,k=1,\ldots,m,
\]
并记 $f_j=e_1-\dfrac{\ii q_0}{\zeta_j}e_2$、$g_j=e_2-\dfrac{\ii q_0}{\conj{\zeta_j}}e_1$. 将 $\mathbf v_j^T,\mathbf u_j^T,f_j^T,g_j^T$ 分别作为 $m\times2$ 矩阵 $\boldsymbol v,\boldsymbol u,\boldsymbol f,\boldsymbol g$ 的第 $j$ 行. 则留数系统为
\begin{equation}\label{eq:N-block-system}
K_N\begin{pmatrix}\boldsymbol v\\\boldsymbol u\end{pmatrix}
=\begin{pmatrix}\boldsymbol f\\\boldsymbol g\end{pmatrix},\qquad
K_N=\begin{pmatrix}I_m&A\\-B&I_m\end{pmatrix}.
\end{equation}
置 $L_N=I_m+AB$. 若 $\det L_N\ne0$, 则系统的唯一解为
\begin{equation}\label{eq:finite-schur}
\boldsymbol v=L_N^{-1}(\boldsymbol f-A\boldsymbol g),\qquad
\boldsymbol u=\boldsymbol g+B\boldsymbol v.
\end{equation}
相应光场为
\begin{equation}\label{eq:qfinite}
q=q_0+\ii\sum\limits_{j=1}^{m}C_j(\mathbf v_j)_1.
\end{equation}
\end{proposition}
\begin{proof}
减去 $z=0$ 处主部及全部离散主部后, 剩余矩阵在有限平面解析, 且在无穷远趋于 $I$. Liouville 定理给出 \eqref{eq:partialfraction}, 这也正是 \eqref{eq:general-cauchy} 中 $G=0$ 的情形. 在第一列取 $z=\zeta_j$, 在第二列取 $z=\conj{\zeta_j}$, 即得 \eqref{eq:finite1}--\eqref{eq:finite2}. 反向代入时, 这些取值等式保证留数系数恰为对应正则列. 写成
\[
\boldsymbol v=\boldsymbol f-A\boldsymbol u,\qquad
\boldsymbol u=\boldsymbol g+B\boldsymbol v,
\]
对系数矩阵作分块消元,
\[
\begin{pmatrix}I_m&-A\\0&I_m\end{pmatrix}
\begin{pmatrix}I_m&A\\-B&I_m\end{pmatrix}
=\begin{pmatrix}L_N&0\\-B&I_m\end{pmatrix},
\qquad \det K_N=\det L_N.
\]
当此行列式非零时, 第一行解出 $\boldsymbol v$, 第二行给出 $\boldsymbol u$, 得到 \eqref{eq:finite-schur}. 最后读取 \eqref{eq:partialfraction} 的 $(1,2)$ 元的 $z^{-1}$ 系数, 得到 \eqref{eq:qfinite}.
\end{proof}

\begin{theorem}[$N$ 轨道矩阵解与行列式重构]\label{thm:N-orbit-RHP}
设数据满足 \eqref{eq:N-orbit-data}--\eqref{eq:N-orbit-phase}. 在任一满足 $\det L_N(x,t)\ne0$ 的实点 $(x,t)$, 式 \eqref{eq:finite-schur} 和 \eqref{eq:partialfraction} 给出无反射 RHP 的唯一解. 它满足 \eqref{eq:Msym} 及 $\det M=\gamma$, 在分支点解析且秩为一. 令
\begin{equation}\label{eq:N-rhs}
\begin{aligned}
&\boldsymbol C=(C_1,\ldots,C_m)^T,\qquad
\boldsymbol r=(\boldsymbol f-A\boldsymbol g)e_1,\\
&r_j=1+\ii q_0\sum\limits_{k=1}^{m}
\dfrac{\conj{C_k}}{(\zeta_j-\conj{\zeta_k})\conj{\zeta_k}},\qquad j=1,\ldots,m,
\end{aligned}
\end{equation}
则规范光场有显式矩阵式与行列式表示
\begin{equation}\label{eq:N-determinant-q}
q=q_0+\ii\boldsymbol C^T L_N^{-1}\boldsymbol r
=\dfrac{\det\begin{pmatrix}L_N&\boldsymbol r\\-\ii\boldsymbol C^T&q_0\end{pmatrix}}
{\det L_N}.
\end{equation}
介质矩阵由同一有理矩阵在共振原像处重构:
\begin{equation}\label{eq:N-medium}
R=M(z_a^\pm)\dfrac{\eta_0\lambda(z_a^\pm)}{a}\sigma_3 M(z_a^\pm)^{-1}.
\end{equation}
两个符号给出相同结果; 原三场按 \eqref{eq:original-fields} 恢复.
\end{theorem}
\begin{proof}
命题 \ref{prop:finite} 给出有理矩阵的全部留数及归一化, 并由 $\det L_N\ne0$ 保证唯一性. 共轭变换 $\sigma_2\conj{M(\conj{z})}\sigma_2$ 交换两组极点, 满足同一留数系统和两处归一化. 换叶变换
\[
\widehat M(z)=M(h(z))\left(\dfrac{\ii z}{q_0}\sigma_1\right)^{-1}
\]
也满足这些条件. 具体地, 记同一轨道的配对指标为 $\widehat j$, 即 $\zeta_{\widehat j}=h(\conj{\zeta_j})$. 在 $z=\zeta_j$ 处, $h(z)$ 趋于 $\conj{\zeta_{\widehat j}}$, 由 $h'(\zeta_j)=\dfrac{q_0^2}{\zeta_j^2}$ 和 \eqref{eq:Cpair},
\[
\Res_{z=\zeta_j}\widehat M_2(z)
=\dfrac{\ii q_0\conj{C_{\widehat j}}}{\zeta_jh'(\zeta_j)}\mathbf u_{\widehat j}
=C_j\left(-\dfrac{\ii q_0}{\zeta_j}\mathbf u_{\widehat j}\right)
=C_j\widehat M_1(\zeta_j).
\]
共轭点处同理, 得到带负号的第一列留数条件. 在 $0,\infty$ 处, 换叶交换两处归一化, 仍给出规定主项. 由唯一性得到两项约化 \eqref{eq:Msym}.

在每个离散极点, 奇异列的留数为正则列的标量倍数, 因而 $\det M$ 的留数为零. 又由 $0,\infty$ 的归一化,
\[
\det M(z)=1+\dfrac{d}{z}+\dfrac{q_0^2}{z^2}
\]
对某个标量 $d$ 成立. 换叶约化的行列式给出
\[
1-\dfrac{dz}{q_0^2}+\dfrac{z^2}{q_0^2}
=\det M(h(z))=\dfrac{z^2}{q_0^2}\det M(z)
=\dfrac{z^2}{q_0^2}+\dfrac{dz}{q_0^2}+1,
\]
故 $d=0$, 即 $\det M=\gamma$. 分支点不在部分分式的极点集合中, 所以 $M$ 在该处解析. 若 $M(z_b)$ 为零矩阵, 则行列式至少有二阶零点, 与 $\gamma'(z_b)\ne0$ 矛盾; 因此其秩恰为一.

由 \eqref{eq:finite-schur}, 正则列第一分量排成的列向量为 $L_N^{-1}\boldsymbol r$. 代入 \eqref{eq:qfinite} 得到 \eqref{eq:N-determinant-q} 的第一式. 再对其分子作分块行列式计算,
\[
\det\begin{pmatrix}L_N&\boldsymbol r\\-\ii\boldsymbol C^T&q_0\end{pmatrix}
=\det L_N\left(q_0+\ii\boldsymbol C^T L_N^{-1}\boldsymbol r\right),
\]
得到第二式. 共振原像不在离散谱中, 且 $\gamma(z_a^\pm)>0$, 所以逆矩阵存在. 两个原像互为换叶, $\lambda(h(z))=-\lambda(z)$, 而 $\sigma_1\sigma_3\sigma_1=-\sigma_3$; 代入换叶约化即得 \eqref{eq:N-medium} 与谱叶无关.
\end{proof}

% BEGIN EXPANSION 08-N
\subsubsection*{完整多轨道数据、Schur 消元与行列式公式的逐项计算}
\paragraph{一个轨道为什么贡献两个正则列指标.}
每个外部上半平面的代表 $\zeta_n$ 对应四点
\[
\zeta_n,\quad \conj{\zeta_n},\quad
h(\zeta_n)=-\dfrac{q_0^2}{\zeta_n},\quad
h(\conj{\zeta_n})=-\dfrac{q_0^2}{\conj{\zeta_n}}.
\]
其中第一点与第四点属于 $D_+$, 第二点与第三点属于 $D_-$.
因此 $N$ 个完整轨道共有 $m=2N$ 个 $D_+$ 点,
而不是只取 $N$ 个外部代表就结束求和.
第 $n$ 与第 $N+n$ 个指标分别指第一点和第四点;
它们的共轭点自动补齐另一个解析域.
因为 $\zeta_n\notin\Sigma$, $\zeta_n\ne0$,
且不同轨道不相交, 所有差 $\zeta_j-\conj{\zeta_k}$ 均不为零.
所以 $A,B$ 的矩阵元在给定数据下都有定义.

由 $\theta(h(\conj{\zeta_n}))=-\conj{\theta(\zeta_n)}$,
\[
\begin{aligned}
C_{N+n}
&=-\dfrac{q_0^2}{(\conj{\zeta_n})^2}\conj{c_n}
e^{-2\ii\conj{\theta(\zeta_n)}}\\
&=-\dfrac{q_0^2}{(\conj{\zeta_n})^2}\conj{C_n}.
\end{aligned}
\]
这说明时空指数保持范数常数的配对关系.
同一轨道对相位乘积的贡献逐因子为
\[
\dfrac{\zeta_n}{\conj{\zeta_n}}
\dfrac{\zeta_{N+n}}{\conj{\zeta_{N+n}}}
=\dfrac{\zeta_n}{\conj{\zeta_n}}
\dfrac{-\dfrac{q_0^2}{\conj{\zeta_n}}}{-\dfrac{q_0^2}{\zeta_n}}
=\left(\dfrac{\zeta_n}{\conj{\zeta_n}}\right)^2.
\]
于是 \eqref{eq:N-orbit-phase} 与上一小节的相位和条件一致.
该约束规定两端背景的相位, 不等同于有限维矩阵的可逆性条件.

\paragraph{把有理矩阵的两列与四个元全部写出.}
由部分分式式 \eqref{eq:partialfraction},
\[
\begin{aligned}
M_1(z)&=e_1-\dfrac{\ii q_0}{z}e_2
-\sum\limits_{k=1}^{m}\dfrac{\conj{C_k}\mathbf u_k}{z-\conj{\zeta_k}},\\
M_2(z)&=e_2-\dfrac{\ii q_0}{z}e_1
+\sum\limits_{k=1}^{m}\dfrac{C_k\mathbf v_k}{z-\zeta_k}.
\end{aligned}
\]
因此
\[
\begin{aligned}
M_{11}(z)&=1-\sum\limits_{k=1}^{m}
\dfrac{\conj{C_k}(\mathbf u_k)_1}{z-\conj{\zeta_k}},\\
M_{21}(z)&=-\dfrac{\ii q_0}{z}-\sum\limits_{k=1}^{m}
\dfrac{\conj{C_k}(\mathbf u_k)_2}{z-\conj{\zeta_k}},\\
M_{12}(z)&=-\dfrac{\ii q_0}{z}+\sum\limits_{k=1}^{m}
\dfrac{C_k(\mathbf v_k)_1}{z-\zeta_k},\\
M_{22}(z)&=1+\sum\limits_{k=1}^{m}
\dfrac{C_k(\mathbf v_k)_2}{z-\zeta_k}.
\end{aligned}
\]
第一列在每个 $\zeta_j$ 处正则, 第二列在每个 $\conj{\zeta_j}$ 处正则.
分别取值就得到 \eqref{eq:finite1}--\eqref{eq:finite2}.
反向代入时, 第二列在 $\zeta_j$ 的留数是 $C_j\mathbf v_j$,
而系统保证 $\mathbf v_j=M_1(\zeta_j)$;
第一列的负号也由相同取值核对, 没有额外的未知留数.

\paragraph{为什么 $4N$ 阶系统能缩减为 $2N$ 阶系统.}
$\boldsymbol v,\boldsymbol u$ 各有 $m$ 行, 每行有两个标量分量.
对任意一个固定的分量, 系数矩阵 $K_N$ 是 $2m\times2m$, 即 $4N$ 阶.
两列右端共用同一系数矩阵, 所以只须分解该矩阵一次.
从第二组方程逐指标代入第一组,
\[
\begin{aligned}
\mathbf v_j
&=f_j-\sum\limits_{k=1}^{m}A_{jk}
\left(g_k+\sum\limits_{l=1}^{m}B_{kl}\mathbf v_l\right)\\
&=f_j-\sum\limits_{k=1}^{m}A_{jk}g_k
-\sum\limits_{l=1}^{m}\left(\sum\limits_{k=1}^{m}A_{jk}B_{kl}\right)\mathbf v_l.
\end{aligned}
\]
把含未知列的项移到左边, 就是
\[
\sum\limits_{l=1}^{m}(L_N)_{jl}\mathbf v_l
=f_j-\sum\limits_{k=1}^{m}A_{jk}g_k,
\]
其中
\[
(L_N)_{jl}=\delta_{jl}+
\sum\limits_{k=1}^{m}
\dfrac{\conj{C_k}C_l}
{(\zeta_j-\conj{\zeta_k})(\conj{\zeta_k}-\zeta_l)}.
\]
这里 $\delta_{jl}$ 为 Kronecker 符号: $j=l$ 时等于 1, 否则等于 0.
这正是矩阵乘积 $L_N=I_m+AB$ 的逐元表达式.
分块消元左乘矩阵的行列式为 1, 因而
$\det K_N=\det L_N$.
所以 $\det L_N\ne0$ 同时保证两个向量分量的唯一可解性,
而不只是其中某个标量方程有解.

右端也可逐分量核对:
\[
\begin{aligned}
(\boldsymbol f-A\boldsymbol g)_{j1}
&=1+\ii q_0\sum\limits_{k=1}^{m}
\dfrac{\conj{C_k}}{(\zeta_j-\conj{\zeta_k})\conj{\zeta_k}}=r_j,\\
(\boldsymbol f-A\boldsymbol g)_{j2}
&=-\dfrac{\ii q_0}{\zeta_j}
-\sum\limits_{k=1}^{m}\dfrac{\conj{C_k}}{\zeta_j-\conj{\zeta_k}}.
\end{aligned}
\]
因此先求 $\boldsymbol v=L_N^{-1}(\boldsymbol f-A\boldsymbol g)$,
再按 $\boldsymbol u=\boldsymbol g+B\boldsymbol v$ 代回,
就得到了全部 $4N$ 条正则列, 并可直接代入上面的四个矩阵元.

\paragraph{光场的边框行列式不省略消元步骤.}
因为 $\boldsymbol v e_1=L_N^{-1}\boldsymbol r$,
\[
\sum\limits_{j=1}^{m}C_j(\mathbf v_j)_1
=\boldsymbol C^T\boldsymbol v e_1
=\boldsymbol C^T L_N^{-1}\boldsymbol r.
\]
此处是转置 $T$, 没有对 $C_j$ 再作共轭.
由 $M_{12}$ 的 $z^{-1}$ 系数乘以 $\ii$, 就得到
$q=q_0+\ii\boldsymbol C^T L_N^{-1}\boldsymbol r$.
再作一次分块消元,
\[
\begin{aligned}
&
\begin{pmatrix}I_m&0\\\ii\boldsymbol C^T L_N^{-1}&1\end{pmatrix}
\begin{pmatrix}L_N&\boldsymbol r\\-\ii\boldsymbol C^T&q_0\end{pmatrix}\\
&\qquad=\begin{pmatrix}L_N&\boldsymbol r\\0&q_0+\ii\boldsymbol C^T L_N^{-1}\boldsymbol r\end{pmatrix}.
\end{aligned}
\]
左乘矩阵行列式仍为 1. 对右边分块三角矩阵取行列式,
再除以非零的 $\det L_N$, 得到 \eqref{eq:N-determinant-q}.
这给出任意有限 $N$ 的直接代数求解, 不需要把单轨道公式反复相加.

\paragraph{换叶留数中的链式因子与负号.}
右乘 $\sigma_1$ 交换两列, 且
$\left(\dfrac{\ii z}{q_0}\sigma_1\right)^{-1}
=-\dfrac{\ii q_0}{z}\sigma_1$.
所以
\[
\widehat M_1(z)=-\dfrac{\ii q_0}{z}M_2(h(z)),\qquad
\widehat M_2(z)=-\dfrac{\ii q_0}{z}M_1(h(z)).
\]
在 $z=\zeta_j$ 附近,
\[
h(z)-\conj{\zeta_{\widehat j}}
=h'(\zeta_j)(z-\zeta_j)+O((z-\zeta_j)^2).
\]
$M_1$ 在共轭点的留数为 $-\conj{C_{\widehat j}}\mathbf u_{\widehat j}$,
与前面的 $-\ii q_0/z$ 相乘给出原证明的正号.
再用
\[
\conj{C_{\widehat j}}=-\dfrac{q_0^2}{\zeta_j^2}C_j,
\qquad h'(\zeta_j)=\dfrac{q_0^2}{\zeta_j^2},
\]
得到原证明中单行等式的每一个因子.
这一步是在完整轨道中配对, 因而对任意 $N$ 都成立.

\paragraph{行列式与分支点的正则性分开检查.}
在 $z=\zeta_j$ 处写第二列为
$\dfrac{C_jM_1(\zeta_j)}{z-\zeta_j}+O(1)$,
则行列式的留数为
\[
C_j\det(M_1(\zeta_j),M_1(\zeta_j))=0.
\]
共轭点同理. 因此行列式只可能在 $0$ 处有二阶极点,
其最高阶系数为
\[
\det\left(-\dfrac{\ii q_0}{z}\sigma_1\right)
=\left(-\dfrac{\ii q_0}{z}\right)^2\det\sigma_1
=\dfrac{q_0^2}{z^2}.
\]
这与无穷远极限 1 一起给出原证明中的 $1+\dfrac{d}{z}+\dfrac{q_0^2}{z^2}$.
换叶后的行列式比较消去 $d$, 从而 $\det M=\gamma$.
在 $z_b=\pm\ii q_0$ 处,
\[
\gamma'(z_b)=-\dfrac{2q_0^2}{z_b^3}\ne0.
\]
这里有一阶零点的是行列式, 不是说矩阵的四个元同时为零.
若四个元都在 $z_b$ 为零, 它们都含因子 $z-z_b$,
行列式便含 $(z-z_b)^2$, 与上述导数矛盾.
因此正则矩阵 $M(z_b)$ 的秩恰为一.

\paragraph{两个共振值以及一般公式的可逆性范围.}
记换叶因子仍为原文的 $\dfrac{\ii z}{q_0}\sigma_1$.
在共振点两叶上, 由 $\lambda(h(z))=-\lambda(z)$,
\[
\begin{aligned}
&M(h(z))\lambda(h(z))\sg M(h(z))^{-1}\\
&\quad=M(z)\sigma_1(-\lambda(z)\sg)\sigma_1M(z)^{-1}
=M(z)\lambda(z)\sg M(z)^{-1}.
\end{aligned}
\]
所以同一部分分式矩阵给出唯一的介质矩阵,
而 $\gamma(z_a^\pm)>0$ 保证这里所用的二阶逆矩阵存在.
计算时先完成有限系统, 再代入两个实共振原像;
无须在共振点计算含极点的传播指数.

一般 $N$ 的公式以 $\det L_N(x,t)\ne0$ 为条件.
它说明在这些点上矩阵解存在且唯一,
并不由相位乘积条件自动推出 $\det L_N$ 在全平面不为零.
下节取 $N=1$ 并指定纯虚数据后,
另行分解行列式并证明严格正性;
该全局正性证明与一般 $N$ 的条件性代数公式应分别阅读.
% END EXPANSION 08-N

上述公式将 $N$ 轨道问题化为 $2N$ 阶矩阵 $L_N$ 的求逆. 公式适用于 $\det L_N\ne0$ 的区域; 下节对一个纯虚轨道进一步显式消元, 并证明该行列式在全体实 $x,t$ 上严格为正.

\subsection{单轨道解及 HMB 方程验证}

% BEGIN EXPANSION 08-E
\paragraph{单轨道的两点系统逐行消元.}
这里取 $N=1$, 因而 $m=2$. 无反射意味着 $G=0$, 因而 $W=0$.
一般表示中连续谱积分逐项等于零, 剩下的不是积分方程, 而是有限个列向量的线性代数方程.
当 $j,k=1,2$ 时, 四行逐行写为
\[
\begin{aligned}
\mathbf v_1
&=f_1- A_{11}\mathbf u_1- A_{12}\mathbf u_2,\\
\mathbf v_2
&=f_2- A_{21}\mathbf u_1- A_{22}\mathbf u_2,\\
\mathbf u_1
&=g_1+ B_{11}\mathbf v_1+ B_{12}\mathbf v_2,\\
\mathbf u_2
&=g_2+ B_{21}\mathbf v_1+ B_{22}\mathbf v_2.
\end{aligned}
\]
把后两行逐个代进第一行, 得
\[
\begin{aligned}
\mathbf v_1
={}&f_1- A_{11}g_1- A_{12}g_2\\
&-( A_{11} B_{11}
+ A_{12} B_{21})\mathbf v_1\\
&-( A_{11} B_{12}
+ A_{12} B_{22})\mathbf v_2.
\end{aligned}
\]
同样代进第二行, 得
\[
\begin{aligned}
\mathbf v_2
={}&f_2- A_{21}g_1- A_{22}g_2\\
&-( A_{21} B_{11}
+ A_{22} B_{21})\mathbf v_1\\
&-( A_{21} B_{12}
+ A_{22} B_{22})\mathbf v_2.
\end{aligned}
\]
把含未知列的项移到左侧, 就是
\[
(I_2+ A B)\boldsymbol v
=\boldsymbol f- A\boldsymbol g.
\]
这里 $I_2+ A B$ 的下标是离散点下标;
每个未知量本身仍是一条二维列.
对每个向量分量分别取值, 所以同一个二阶系数矩阵解两次.
这也解释原文将 $\boldsymbol v$ 记成两行两列数组时, 每一行实际是一条正则列的转置.
% END EXPANSION 08-E
取一个纯虚离散轨道
\begin{equation}\label{eq:imaginarypoint}
k_1=\ii b,\qquad b>q_0,\qquad
\kappa=\sqrt{b^2-q_0^2},\qquad W=b+\kappa,\qquad h_0=\dfrac{q_0}{W}\in(0,1).
\end{equation}
两个 $D_+$ 零点与透射函数为
\begin{equation}\label{eq:onezeros}
\zeta_1=\ii W,\qquad \zeta_2=-\ii q_0h_0=h(\conj{\zeta_1}),\qquad
s_{22}(z)=\dfrac{\lambda(z)-\ii\kappa}{\lambda(z)+\ii\kappa}.
\end{equation}
这正是无反射迹公式的两因子乘积, 其在 $0,\infty$ 均趋于 1. 在外部零点选择 $b_1=1$, 并由换叶约化确定另一常数:
\begin{equation}\label{eq:onec}
c_1=\dfrac{1}{s_{22}'(\zeta_1)}=\dfrac{4\ii\kappa}{1+h_0^2},\qquad
c_2=-h_0^2c_1.
\end{equation}
令
\begin{equation}\label{eq:orbit-phase}
s=\kappa[x-F(\ii b)t],\qquad T=e^{2s}.
\end{equation}
由 $\lambda(\zeta_1)=\lambda(\zeta_2)=\ii\kappa$, $k(\zeta_2)=-\ii b$ 及 $F(-\ii b)=\conj{F(\ii b)}$, 得 $C_1=c_1T$、$C_2=-h_0^2c_1\conj{T}$.

% BEGIN EXPANSION 08-F
\paragraph{谱点, 换叶, 透射导数与四个范数系数的实算.}
由 $\kappa^2=b^2-q_0^2$ 得
\[
(b+\kappa)(b-\kappa)=q_0^2.
\]
结合 $W=b+\kappa$ 与 $h_0=\dfrac{q_0}{W}$, 逐式得到
\[
\begin{aligned}
W&=\dfrac{q_0}{h_0},&
b-\kappa&=q_0h_0,\\
2b&=\dfrac{q_0}{h_0}+q_0h_0
=\dfrac{q_0(1+h_0^2)}{h_0},\\
2\kappa&=\dfrac{q_0}{h_0}-q_0h_0
=\dfrac{q_0(1-h_0^2)}{h_0}.
\end{aligned}
\]
由于 $b>q_0>0$, 有 $\kappa>0$ 和 $W>q_0$, 所以 $0<h_0<1$.
四个点逐一列出是
\[
\zeta_1=\dfrac{\ii q_0}{h_0},\quad
\zeta_2=-\ii q_0h_0,\quad
\conj{\zeta_1}=-\dfrac{\ii q_0}{h_0},\quad
\conj{\zeta_2}=\ii q_0h_0.
\]
在第一个点代入均匀化公式:
\[
\begin{aligned}
\lambda(\zeta_1)
&=\dfrac{1}{2}\left(\dfrac{\ii q_0}{h_0}
+\dfrac{q_0^2}{\dfrac{\ii q_0}{h_0}}\right)
=\dfrac{\ii q_0}{2}(h_0^{-1}-h_0)=\ii\kappa,\\
k(\zeta_1)
&=\dfrac{\ii q_0}{2}(h_0^{-1}+h_0)=\ii b.
\end{aligned}
\]
在第二个点同样计算:
\[
\begin{aligned}
\lambda(\zeta_2)
&=\dfrac{1}{2}\left(-\ii q_0h_0
+\dfrac{q_0^2}{-\ii q_0h_0}\right)
=\dfrac{\ii q_0}{2}(h_0^{-1}-h_0)=\ii\kappa,\\
k(\zeta_2)
&=-\dfrac{\ii q_0}{2}(h_0+h_0^{-1})=-\ii b.
\end{aligned}
\]
共轭点的对应值则是
\[
\begin{aligned}
\lambda(\conj{\zeta_1})&=-\ii\kappa,&
k(\conj{\zeta_1})&=-\ii b,\\
\lambda(\conj{\zeta_2})&=-\ii\kappa,&
k(\conj{\zeta_2})&=\ii b.
\end{aligned}
\]
两个零因子的分子相乘为
\[
\begin{aligned}
(z-\zeta_1)(z-\zeta_2)
&=\left(z-\dfrac{\ii q_0}{h_0}\right)
\left(z+\ii q_0h_0\right)\\
&=z^2-\ii q_0(h_0^{-1}-h_0)z+q_0^2\\
&=z^2-2\ii\kappa z+q_0^2
=2z(\lambda-\ii\kappa).
\end{aligned}
\]
共轭分母相乘同理为 $2z(\lambda+\ii\kappa)$.
两者相除便是原文的 $s_{22}$.
对该式按 $z$ 求导, 使用商法则:
\[
\begin{aligned}
s_{22}'(z)
&=\dfrac{\lambda'(z)(\lambda+\ii\kappa)
-(\lambda-\ii\kappa)\lambda'(z)}
{(\lambda+\ii\kappa)^2}\\
&=\dfrac{2\ii\kappa\lambda'(z)}
{(\lambda+\ii\kappa)^2},\\
\lambda'(z)&=\dfrac{1}{2}\left(1-\dfrac{q_0^2}{z^2}\right).
\end{aligned}
\]
四点上的导数分别为
\[
\begin{aligned}
\lambda'(\zeta_1)=\lambda'(\conj{\zeta_1})
&=\dfrac{1+h_0^2}{2},\\
\lambda'(\zeta_2)=\lambda'(\conj{\zeta_2})
&=\dfrac{1+h_0^{-2}}{2}
=\dfrac{1+h_0^2}{2h_0^2}.
\end{aligned}
\]
所以
\[
\begin{aligned}
s_{22}'(\zeta_1)
&=\dfrac{\lambda'(\zeta_1)}{2\ii\kappa}
=\dfrac{1+h_0^2}{4\ii\kappa},\\
s_{22}'(\zeta_2)
&=\dfrac{1+h_0^2}{4\ii\kappa h_0^2}.
\end{aligned}
\]
原选择 $b_1=1$ 与约化 $b_2=-\conj{b_1}=-1$ 给出
\[
\begin{aligned}
c_1&=\dfrac{1}{s_{22}'(\zeta_1)}
=\dfrac{4\ii\kappa}{1+h_0^2},\\
c_2&=\dfrac{-1}{s_{22}'(\zeta_2)}
=-\dfrac{4\ii\kappa h_0^2}{1+h_0^2}
=-h_0^2c_1.
\end{aligned}
\]
因此 $\conj{c_1}=-c_1$, $\conj{c_2}=-c_2$.
再计算时间指数:
\[
\begin{aligned}
2\ii\theta(\zeta_1)
&=2\ii(-\ii\kappa)[x-F(\ii b)t]
=2\kappa[x-F(\ii b)t]=2s,\\
2\ii\theta(\zeta_2)
&=2\kappa[x-F(-\ii b)t]=2\conj s.
\end{aligned}
\]
所以四个实际留数标量是
\[
\begin{aligned}
C_1&=\dfrac{4\ii\kappa}{1+h_0^2}T,&
C_2&=-\dfrac{4\ii\kappa h_0^2}{1+h_0^2}\conj T,\\
\conj{C_1}&=-\dfrac{4\ii\kappa}{1+h_0^2}\conj T,&
\conj{C_2}&=\dfrac{4\ii\kappa h_0^2}{1+h_0^2}T.
\end{aligned}
\]
% END EXPANSION 08-F
\begin{lemma}[一轨道系统的消元]\label{lem:direct-solve}
上述留数系统对任意实 $x,t$ 唯一可解. 记
\begin{equation}\label{eq:delta-pair}
\Delta_\pm=(1+h_0^2)(1+|T|^2)\pm2h_0(T+\conj{T}),
\end{equation}
则其解为
\begin{align}
\mathbf v_1&=\dfrac{1+h_0^2}{\Delta_-}
\begin{pmatrix}1-h_0\conj{T}\\\conj{T}-h_0\end{pmatrix},&
\mathbf u_1&=\dfrac{1+h_0^2}{\Delta_-}
\begin{pmatrix}h_0-T\\1-h_0T\end{pmatrix},\label{eq:direct-vu1}\\
\mathbf v_2&=\dfrac{1}{h_0}\mathbf u_1,&
\mathbf u_2&=-\dfrac{1}{h_0}\mathbf v_1.\label{eq:direct-vu2}
\end{align}
\end{lemma}
\begin{proof}
代入谱点和范数常数得到
\begin{equation}\label{eq:one-AB}
A=\dfrac{1}{1+h_0^2}
\begin{pmatrix}(h_0^2-1)\conj{T}&2h_0^2T\\-2\conj{T}&(h_0^2-1)T\end{pmatrix},\quad
B=\dfrac{1}{1+h_0^2}
\begin{pmatrix}(h_0^2-1)T&2h_0^2\conj{T}\\-2T&(h_0^2-1)\conj{T}\end{pmatrix}.
\end{equation}
对应右端为
\[
f_1=\begin{pmatrix}1\\-h_0\end{pmatrix},\quad
f_2=\begin{pmatrix}1\\h_0^{-1}\end{pmatrix},\quad
g_1=\begin{pmatrix}h_0\\1\end{pmatrix},\quad
g_2=\begin{pmatrix}-h_0^{-1}\\1\end{pmatrix}.
\]
两个向量分量分别满足同一个四阶标量系数系统,
\[
\begin{pmatrix}I_2&A\\-B&I_2\end{pmatrix}
\begin{pmatrix}\boldsymbol v\\\boldsymbol u\end{pmatrix}
=\begin{pmatrix}\boldsymbol f\\\boldsymbol g\end{pmatrix}.
\]
其 Schur 补行列式为
\begin{equation}\label{eq:one-schur-det}
\det(I_2+AB)=\dfrac{\Delta_-\Delta_+}{(1+h_0^2)^2}>0.
\end{equation}
事实上,
\[
\Delta_-=|T-h_0|^2+|1-h_0T|^2,\qquad
\Delta_+=|T+h_0|^2+|1+h_0T|^2,
\]
且 $\Delta_\pm\ge(1-h_0)^2(1+|T|^2)$. 令 $L=I_2+AB$, 用二阶逆矩阵公式计算 $L^{-1}(\boldsymbol f-A\boldsymbol g)$, 约去公因子 $\Delta_+$, 得
\[
\boldsymbol v=\dfrac{1+h_0^2}{\Delta_-}
\begin{pmatrix}1-h_0\conj{T}&\conj{T}-h_0\\
\dfrac{h_0-T}{h_0}&\dfrac{1-h_0T}{h_0}\end{pmatrix}.
\]
此处每一行是一个正则列的转置. 由 $\boldsymbol u=\boldsymbol g+B\boldsymbol v$ 得到其余两列, 即 \eqref{eq:direct-vu1}--\eqref{eq:direct-vu2}.
\end{proof}

% BEGIN EXPANSION 08-G
\subsubsection*{留数系统的八个系数, Schur 补的四元与二阶逆矩阵}
四种差分母逐个为
\[
\begin{aligned}
\zeta_1-\conj{\zeta_1}&=\dfrac{2\ii q_0}{h_0},\\
\zeta_1-\conj{\zeta_2}&=\dfrac{\ii q_0(1-h_0^2)}{h_0},\\
\zeta_2-\conj{\zeta_1}&=\dfrac{\ii q_0(1-h_0^2)}{h_0},\\
\zeta_2-\conj{\zeta_2}&=-2\ii q_0h_0.
\end{aligned}
\]
由 $2\kappa h_0=q_0(1-h_0^2)$, 八个商逐个化简为
\[
\begin{aligned}
 A_{11}
&=\dfrac{-4\ii\kappa\conj T}{1+h_0^2}
\dfrac{h_0}{2\ii q_0}
=\dfrac{(h_0^2-1)\conj T}{1+h_0^2},\\
 A_{12}
&=\dfrac{4\ii\kappa h_0^2T}{1+h_0^2}
\dfrac{h_0}{\ii q_0(1-h_0^2)}
=\dfrac{2h_0^2T}{1+h_0^2},\\
 A_{21}
&=\dfrac{-4\ii\kappa\conj T}{1+h_0^2}
\dfrac{h_0}{\ii q_0(1-h_0^2)}
=-\dfrac{2\conj T}{1+h_0^2},\\
 A_{22}
&=\dfrac{4\ii\kappa h_0^2T}{1+h_0^2}
\dfrac{-1}{2\ii q_0h_0}
=\dfrac{(h_0^2-1)T}{1+h_0^2},\\
 B_{11}
&=\dfrac{4\ii\kappa T}{1+h_0^2}
\dfrac{-h_0}{2\ii q_0}
=\dfrac{(h_0^2-1)T}{1+h_0^2},\\
 B_{12}
&=\dfrac{-4\ii\kappa h_0^2\conj T}{1+h_0^2}
\dfrac{-h_0}{\ii q_0(1-h_0^2)}
=\dfrac{2h_0^2\conj T}{1+h_0^2},\\
 B_{21}
&=\dfrac{4\ii\kappa T}{1+h_0^2}
\dfrac{-h_0}{\ii q_0(1-h_0^2)}
=-\dfrac{2T}{1+h_0^2},\\
 B_{22}
&=\dfrac{-4\ii\kappa h_0^2\conj T}{1+h_0^2}
\dfrac{1}{2\ii q_0h_0}
=\dfrac{(h_0^2-1)\conj T}{1+h_0^2}.
\end{aligned}
\]
原文已经定义 $L=I_2+ A B$.
矩阵乘法每个元只包含两个乘积, 全部写为
\[
\begin{aligned}
L_{11}
&=1+ A_{11} B_{11}
+ A_{12} B_{21}\\
&=1+\dfrac{(h_0^2-1)^2|T|^2-4h_0^2T^2}
{(1+h_0^2)^2},\\
L_{12}
&= A_{11} B_{12}
+ A_{12} B_{22}\\
&=\dfrac{2h_0^2(h_0^2-1)((\conj{T})^2+|T|^2)}
{(1+h_0^2)^2},\\
L_{21}
&= A_{21} B_{11}
+ A_{22} B_{21}\\
&=-\dfrac{2(h_0^2-1)(|T|^2+T^2)}
{(1+h_0^2)^2},\\
L_{22}
&=1+ A_{21} B_{12}
+ A_{22} B_{22}\\
&=1+\dfrac{(h_0^2-1)^2|T|^2-4h_0^2(\conj{T})^2}
{(1+h_0^2)^2}.
\end{aligned}
\]
先分别计算两个系数矩阵的行列式:
\[
\begin{aligned}
\det A
&=\dfrac{(h_0^2-1)^2\conj T T+4h_0^2T\conj T}
{(1+h_0^2)^2}\\
&=\dfrac{(h_0^4-2h_0^2+1+4h_0^2)|T|^2}
{(1+h_0^2)^2}=|T|^2,\\
\det B&=|T|^2.
\end{aligned}
\]
任意二阶矩阵的恒等式
$\det(I_2+ A B)
=1+\tr( A B)+\det( A B)$
可由 $L_{11}L_{22}-L_{12}L_{21}$ 直接展开得到.
这里代入上面的四元, 给出
\[
\det L
=1+|T|^4
+\dfrac{2(h_0^2-1)^2|T|^2
-4h_0^2(T^2+(\conj{T})^2)}{(1+h_0^2)^2}.
\]
另一方面, 把 $\Delta_\pm$ 的乘积按平方差展开:
\[
\begin{aligned}
\Delta_-\Delta_+
={}&(1+h_0^2)^2(1+|T|^2)^2
-4h_0^2(T+\conj T)^2\\
={}&(1+h_0^2)^2(1+|T|^4)\\
&+\bigl(2(1+h_0^2)^2-8h_0^2\bigr)|T|^2
-4h_0^2(T^2+(\conj{T})^2)\\
={}&(1+h_0^2)^2(1+|T|^4)
+2(h_0^2-1)^2|T|^2
-4h_0^2(T^2+(\conj{T})^2).
\end{aligned}
\]
除以 $(1+h_0^2)^2$ 恰为 $\det L$.
正性可不用任何关于 $T$ 相位的限制:
\[
\begin{aligned}
|T\mp h_0|^2+|1\mp h_0T|^2
&=(1+h_0^2)(1+|T|^2)
\mp2h_0(T+\conj T),\\
\Delta_\pm
&\ge(1+h_0^2)(1+|T|^2)-4h_0|T|\\
&=(1-h_0)^2(1+|T|^2)
+2h_0(|T|-1)^2\\
&\ge(1-h_0)^2(1+|T|^2)>0.
\end{aligned}
\]
因此四阶系数系统对于任意复数 $T$ 都可逆;
实际 $T=e^{2s}$ 当然也在这个范围内.

\paragraph{右端四个标量以及伴随矩阵乘法.}
将两条列的转置排成行, 原有右端的四个元分别为
\[
\begin{aligned}
(\boldsymbol f- A\boldsymbol g)_{11}
&=1-h_0 A_{11}+h_0^{-1} A_{12}\\
&=\dfrac{1+h_0^2-h_0(h_0^2-1)\conj T+2h_0T}
{1+h_0^2},\\
(\boldsymbol f- A\boldsymbol g)_{12}
&=-h_0- A_{11}- A_{12}\\
&=\dfrac{-h_0(1+h_0^2)-(h_0^2-1)\conj T-2h_0^2T}
{1+h_0^2},\\
(\boldsymbol f- A\boldsymbol g)_{21}
&=1-h_0 A_{21}+h_0^{-1} A_{22}\\
&=\dfrac{h_0(1+h_0^2)+2h_0^2\conj T+(h_0^2-1)T}
{h_0(1+h_0^2)},\\
(\boldsymbol f- A\boldsymbol g)_{22}
&=h_0^{-1}- A_{21}- A_{22}\\
&=\dfrac{1+h_0^2+2h_0\conj T-h_0(h_0^2-1)T}
{h_0(1+h_0^2)}.
\end{aligned}
\]
二阶逆矩阵的完整式为
\[
L^{-1}
=\dfrac{1}{\det L}
\begin{pmatrix}L_{22}&-L_{12}\\-L_{21}&L_{11}\end{pmatrix}.
\]
例如其与第一列相乘的第一元, 在通分后是
\[
\begin{aligned}
&L_{22}(\boldsymbol f- A\boldsymbol g)_{11}
-L_{12}(\boldsymbol f- A\boldsymbol g)_{21}\\
={}&\dfrac{1}{(1+h_0^2)^3}
\Bigl[
\bigl((1+h_0^2)^2+(h_0^2-1)^2|T|^2-4h_0^2(\conj{T})^2\bigr)\\
&\qquad\cdot
\bigl(1+h_0^2-h_0(h_0^2-1)\conj T+2h_0T\bigr)\\
&\qquad-2h_0(h_0^2-1)((\conj{T})^2+|T|^2)
\bigl(h_0(1+h_0^2)+2h_0^2\conj T+(h_0^2-1)T\bigr)
\Bigr]\\
={}&\dfrac{\Delta_+}{1+h_0^2}(1-h_0\conj T).
\end{aligned}
\]
最后一行可以逐幂核验: 将 $|T|^2$ 写回 $T\conj T$,
同类项相加, 分子正好为
\[
(1+h_0^2)^2(1-h_0\conj T)
\bigl((1+h_0^2)(1+T\conj T)+2h_0(T+\conj T)\bigr).
\]
同一伴随矩阵乘法的其余三个元, 不省略其组合方式:
\[
\begin{aligned}
&L_{22}(\boldsymbol f- A\boldsymbol g)_{12}
-L_{12}(\boldsymbol f- A\boldsymbol g)_{22}
=\dfrac{\Delta_+}{1+h_0^2}(\conj T-h_0),\\
&-L_{21}(\boldsymbol f- A\boldsymbol g)_{11}
+L_{11}(\boldsymbol f- A\boldsymbol g)_{21}
=\dfrac{\Delta_+}{h_0(1+h_0^2)}(h_0-T),\\
&-L_{21}(\boldsymbol f- A\boldsymbol g)_{12}
+L_{11}(\boldsymbol f- A\boldsymbol g)_{22}
=\dfrac{\Delta_+}{h_0(1+h_0^2)}(1-h_0T).
\end{aligned}
\]
用 $\det L=\dfrac{\Delta_-\Delta_+}{(1+h_0^2)^2}$ 除这四项,
每一项中的 $\Delta_+$ 都约掉, 正好得到原文的 $\boldsymbol v$.

\paragraph{将 $\boldsymbol v$ 代回以得到 $\boldsymbol u$.}
例如第一条列的第一元是
\[
\begin{aligned}
(\mathbf u_1)_1
&=h_0+ B_{11}(\mathbf v_1)_1
+ B_{12}(\mathbf v_2)_1\\
&=h_0+\dfrac{(h_0^2-1)T(1-h_0\conj T)
+2h_0\conj T(h_0-T)}{\Delta_-}\\
&=\dfrac{h_0\Delta_-+(h_0^2-1)T(1-h_0\conj T)
+2h_0\conj T(h_0-T)}{\Delta_-}\\
&=\dfrac{(1+h_0^2)(h_0-T)}{\Delta_-}.
\end{aligned}
\]
其余三个元使用同一个代入, 分子分别为
\[
\begin{aligned}
\Delta_-(\mathbf u_1)_2
&=\Delta_-+(h_0^2-1)T(\conj T-h_0)
+2h_0\conj T(1-h_0T)\\
&=(1+h_0^2)(1-h_0T),\\
\Delta_-(\mathbf u_2)_1
&=-\dfrac{\Delta_-}{h_0}-2T(1-h_0\conj T)
+\dfrac{h_0^2-1}{h_0}\conj T(h_0-T)\\
&=-\dfrac{1+h_0^2}{h_0}(1-h_0\conj T),\\
\Delta_-(\mathbf u_2)_2
&=\Delta_--2T(\conj T-h_0)
+\dfrac{h_0^2-1}{h_0}\conj T(1-h_0T)\\
&=-\dfrac{1+h_0^2}{h_0}(\conj T-h_0).
\end{aligned}
\]
这四式既给出 $\mathbf u_1,\mathbf u_2$, 也直接检查
$\mathbf v_2=\dfrac{\mathbf u_1}{h_0}$ 与 $\mathbf u_2=-\dfrac{\mathbf v_1}{h_0}$.
没有使用预先假定的换叶关系去缩小系统;
换叶关系是解出完整系统后得到的结论.
% END EXPANSION 08-G
将所得列向量代入部分分式表示. 其系数可用以下秩一正交投影写成紧凑形式:
\begin{equation}\label{eq:algebraic-P}
P=\dfrac{1}{\Delta_-}
\begin{pmatrix}
|T-h_0|^2&(T-h_0)(1-h_0\conj{T})\\
(1-h_0T)(\conj{T}-h_0)&|1-h_0T|^2
\end{pmatrix},\qquad H=\diag(\ii,-\ii).
\end{equation}
由 $\Delta_-=|T-h_0|^2+|1-h_0T|^2$, 有 $P^\dagger=P$, $P^2=P$, $\tr P=1$.

\begin{theorem}[完整四点轨道的矩阵解]\label{thm:oneRHP}\label{prop:unique-reflectionless}
在背景参数 $q_0>0$, $a\ne0$, $\eta_0\ne0$, $\alpha,\beta\in\mathbb R$ 和 $b>q_0$ 下, 数据 \eqref{eq:onezeros}--\eqref{eq:onec} 的无反射 RHP 有唯一解
\begin{equation}\label{eq:explicitM}
M(z)=H\bigl[(k(z)+\ii b)I-2\ii bP\bigr]Y(z)
\diag\left(\dfrac{1}{\lambda(z)+\ii\kappa},\dfrac{1}{\lambda(z)-\ii\kappa}\right)H^{-1}.
\end{equation}
它满足定理 \ref{thm:RHP} 的全部条件, 在分支点、轮廓交点与共振原像解析, 在分支点秩为一, 在共振原像可逆. 唯一性取于除 $0$ 和指定四点外没有有限奇点、非极点轮廓处可去并满足留数及两处归一化的亚纯矩阵类.
\end{theorem}
\begin{proof}
在部分分式的第二列合并分母,
$(z-\zeta_1)(z-\zeta_2)=2z(\lambda-\ii\kappa)$, 第一列同理. 利用
$2bh_0=q_0(1+h_0^2)$、$2\kappa h_0=q_0(1-h_0^2)$ 化简分子, 得到 \eqref{eq:explicitM}. 留数、两处归一化及恒等跳跃已由命题 \ref{prop:finite} 保证.

共轭矩阵 $\sigma_2\conj{M(\conj{z})}\sigma_2$ 满足同一留数系统, 其中交换两组极点产生 \eqref{eq:resminus} 的负号. 换叶矩阵 $M(h(z))C(z)^{-1}$ 也满足该系统: 留数变换的导数因子由 \eqref{eq:cpair} 抵消, 而 $\theta(h(\conj{\zeta}))=-\conj{\theta(\zeta)}$ 保证指数一致. 两处归一化相互交换后仍为规定主项. 引理 \ref{lem:direct-solve} 的唯一性遂给出 \eqref{eq:Msym}.

记
\begin{equation}\label{eq:Nfactor}
N(k)=(k+\ii b)I-2\ii bP.
\end{equation}
$P$ 的特征值为 $1,0$, 故
\[
\det N=(k-\ii b)(k+\ii b)=\lambda^2+\kappa^2.
\]
它抵消两个列分母的乘积, 得 $\det M=\det Y=\gamma$. 在分支点 $z=\pm\ii q_0$ 上, $\lambda=0$, 两个列分母均非零; $N$ 可逆而 $Y$ 秩一, 因此 $M$ 解析且秩一. 交点和共振原像的 $\lambda$ 为实数, 也没有列分母零点; 共振处 $\gamma>0$, 所以 $M$ 可逆. 显式表达式不再含有其他有限奇点.

任一属于所述函数类的解均具有 \eqref{eq:partialfraction} 的表示, 因而由同一留数系统确定. Schur 行列式 \eqref{eq:one-schur-det} 严格为正, 保证该解唯一.
\end{proof}

% BEGIN EXPANSION 08-H
\subsubsection*{投影, 有理矩阵的四元以及四点留数的直接核验}
先逐项检查原有投影矩阵.
其四元是
\[
\begin{aligned}
P_{11}&=\dfrac{(T-h_0)(\conj T-h_0)}{\Delta_-},\\
P_{12}&=\dfrac{(T-h_0)(1-h_0\conj T)}{\Delta_-},\\
P_{21}&=\dfrac{(1-h_0T)(\conj T-h_0)}{\Delta_-},\\
P_{22}&=\dfrac{(1-h_0T)(1-h_0\conj T)}{\Delta_-}.
\end{aligned}
\]
分母为正实数, 所以 $P_{11},P_{22}$ 为实数且
$\conj{P_{12}}=P_{21}$.
这逐元给出 $P^\dagger=P$.
迹的计算为
\[
P_{11}+P_{22}
=\dfrac{|T-h_0|^2+|1-h_0T|^2}{\Delta_-}=1.
\]
行列式分子中两项完全相同:
\[
\begin{aligned}
\Delta_-^2\det P
={}&(T-h_0)(\conj T-h_0)(1-h_0T)(1-h_0\conj T)\\
&-(T-h_0)(1-h_0\conj T)(1-h_0T)(\conj T-h_0)=0.
\end{aligned}
\]
平方矩阵四元的乘法则是
\[
\begin{aligned}
(P^2)_{11}
&=\dfrac{|T-h_0|^4+|T-h_0|^2|1-h_0T|^2}{\Delta_-^2}
=P_{11},\\
(P^2)_{12}
&=\dfrac{(T-h_0)(1-h_0\conj T)
(|T-h_0|^2+|1-h_0T|^2)}{\Delta_-^2}
=P_{12},\\
(P^2)_{21}
&=\dfrac{(1-h_0T)(\conj T-h_0)
(|T-h_0|^2+|1-h_0T|^2)}{\Delta_-^2}
=P_{21},\\
(P^2)_{22}
&=\dfrac{|1-h_0T|^2|T-h_0|^2+|1-h_0T|^4}{\Delta_-^2}
=P_{22}.
\end{aligned}
\]
其中第二、第三行分子中紧邻括号表示乘积.
因此该矩阵确为正交投影.
其迹为 1 而非零, 行列式为零, 所以秩恰为一.

\paragraph{由部分分式合并得到显式矩阵, 不跳过四个元.}
四个留数列的系数分别是
\[
\begin{aligned}
C_1\mathbf v_1
&=\dfrac{4\ii\kappa T}{\Delta_-}
\begin{pmatrix}1-h_0\conj T\\\conj T-h_0\end{pmatrix},\\
C_2\mathbf v_2
&=\dfrac{4\ii\kappa h_0\conj T}{\Delta_-}
\begin{pmatrix}T-h_0\\h_0T-1\end{pmatrix},\\
-\conj{C_1}\mathbf u_1
&=\dfrac{4\ii\kappa\conj T}{\Delta_-}
\begin{pmatrix}h_0-T\\1-h_0T\end{pmatrix},\\
-\conj{C_2}\mathbf u_2
&=\dfrac{4\ii\kappa h_0T}{\Delta_-}
\begin{pmatrix}1-h_0\conj T\\\conj T-h_0\end{pmatrix}.
\end{aligned}
\]
因此原部分分式四元完整写为
\[
\begin{aligned}
M_{11}
&=1+\dfrac{4\ii\kappa}{\Delta_-}
\left[
\dfrac{\conj T(h_0-T)}{z+\dfrac{\ii q_0}{h_0}}
+\dfrac{h_0T(1-h_0\conj T)}{z-\ii q_0h_0}
\right],\\
M_{21}
&=-\dfrac{\ii q_0}{z}
+\dfrac{4\ii\kappa}{\Delta_-}
\left[
\dfrac{\conj T(1-h_0T)}{z+\dfrac{\ii q_0}{h_0}}
+\dfrac{h_0T(\conj T-h_0)}{z-\ii q_0h_0}
\right],\\
M_{12}
&=-\dfrac{\ii q_0}{z}
+\dfrac{4\ii\kappa}{\Delta_-}
\left[
\dfrac{T(1-h_0\conj T)}{z-\dfrac{\ii q_0}{h_0}}
+\dfrac{h_0\conj T(T-h_0)}{z+\ii q_0h_0}
\right],\\
M_{22}
&=1+\dfrac{4\ii\kappa}{\Delta_-}
\left[
\dfrac{T(\conj T-h_0)}{z-\dfrac{\ii q_0}{h_0}}
-\dfrac{h_0\conj T(1-h_0T)}{z+\ii q_0h_0}
\right].
\end{aligned}
\]
第一列两个分母的乘积是 $2z(\lambda+\ii\kappa)$,
第二列两个分母的乘积是 $2z(\lambda-\ii\kappa)$.
例如 $M_{12}$ 中两个留数项通分时的分子为
\[
\begin{aligned}
&T(1-h_0\conj T)(z+\ii q_0h_0)
+h_0\conj T(T-h_0)(z-\dfrac{\ii q_0}{h_0})\\
={}&z\bigl(T-h_0^2\conj T\bigr)
+\ii q_0\bigl(h_0(T+\conj T)-(1+h_0^2)|T|^2\bigr).
\end{aligned}
\]
其余三个元的对应分子也逐个为
\[
\begin{aligned}
&\conj T(h_0-T)(z-\ii q_0h_0)
+h_0T(1-h_0\conj T)(z+\dfrac{\ii q_0}{h_0})\\
&\qquad=z\bigl(h_0(T+\conj T)-(1+h_0^2)|T|^2\bigr)
+\ii q_0(T-h_0^2\conj T),\\
&\conj T(1-h_0T)(z-\ii q_0h_0)
+h_0T(\conj T-h_0)(z+\dfrac{\ii q_0}{h_0})\\
&\qquad=z(\conj T-h_0^2T)
+\ii q_0\bigl((1+h_0^2)|T|^2-h_0(T+\conj T)\bigr),\\
&T(\conj T-h_0)(z+\ii q_0h_0)
-h_0\conj T(1-h_0T)(z-\dfrac{\ii q_0}{h_0})\\
&\qquad=z\bigl((1+h_0^2)|T|^2-h_0(T+\conj T)\bigr)
+\ii q_0(\conj T-h_0^2T).
\end{aligned}
\]
接着将 $N$ 的四元直接写出:
\[
\begin{aligned}
N_{11}&=k+\ii b-2\ii bP_{11},&
N_{12}&=-2\ii bP_{12},\\
N_{21}&=-2\ii bP_{21},&
N_{22}&=k+\ii b-2\ii bP_{22}.
\end{aligned}
\]
因为
\[
H=\begin{pmatrix}\ii&0\\0&-\ii\end{pmatrix},\qquad
H^{-1}=\begin{pmatrix}-\ii&0\\0&\ii\end{pmatrix},
\]
对任意矩阵左右乘 $H,H^{-1}$ 会保持对角元, 将两个非对角元变号.
因此显式矩阵 \eqref{eq:explicitM} 按元展开为
\[
\begin{aligned}
M_{11}
&=\dfrac{N_{11}-\dfrac{\ii q_0}{z}N_{12}}{\lambda+\ii\kappa},\\
M_{21}
&=-\dfrac{N_{21}-\dfrac{\ii q_0}{z}N_{22}}{\lambda+\ii\kappa},\\
M_{12}
&=-\dfrac{N_{12}-\dfrac{\ii q_0}{z}N_{11}}{\lambda-\ii\kappa},\\
M_{22}
&=\dfrac{N_{22}-\dfrac{\ii q_0}{z}N_{21}}{\lambda-\ii\kappa}.
\end{aligned}
\]
代入 $P_{jk}$ 后, 这四式与上面的通分式逐元相同.
核对时仅需使用
\[
\lambda-k=\dfrac{q_0^2}{z},\qquad
2bh_0=q_0(1+h_0^2),\qquad
2\kappa h_0=q_0(1-h_0^2).
\]
例如第三式的分子可无省略地写为
\[
\begin{aligned}
&2\ii b(T-h_0)(1-h_0\conj T)
+\dfrac{\ii q_0}{z}
\bigl((k+\ii b)\Delta_--2\ii b|T-h_0|^2\bigr)\\
={}&-\ii q_0(\lambda-\ii\kappa)\Delta_-\\
&+\dfrac{2\ii\kappa}{z}
\left[
z(T-h_0^2\conj T)
+\ii q_0\bigl(h_0(T+\conj T)-(1+h_0^2)|T|^2\bigr)
\right].
\end{aligned}
\]
两侧同除以 $(\lambda-\ii\kappa)\Delta_-$ 就得到先前的 $M_{12}$.
这给出了从部分分式到紧凑矩阵乘积的可直接复算过程.

\paragraph{外部 $D_+$ 点 $\zeta_1$ 的正则列与留数.}
在此点 $k=\ii b$, 所以 $N=2\ii b(I-P)$.
两列背景向量分别为 $(1,-h_0)^T,(-h_0,1)^T$.
投影的两次乘法逐项为
\[
\begin{aligned}
(I-P)\begin{pmatrix}1\\-h_0\end{pmatrix}
&=\dfrac{1-h_0^2}{\Delta_-}
\begin{pmatrix}1-h_0\conj T\\-(\conj T-h_0)\end{pmatrix},\\
(I-P)\begin{pmatrix}-h_0\\1\end{pmatrix}
&=\dfrac{(1-h_0^2)T}{\Delta_-}
\begin{pmatrix}-(1-h_0\conj T)\\\conj T-h_0\end{pmatrix}.
\end{aligned}
\]
例如第一式第一元的分子是
\[
|1-h_0T|^2+h_0(T-h_0)(1-h_0\conj T)
=(1-h_0\conj T)(1-h_0^2).
\]
先读正则第一列, 分母是 $2\ii\kappa$:
\[
M_1(\zeta_1)
=\dfrac{b(1-h_0^2)}{\kappa\Delta_-}
\begin{pmatrix}1-h_0\conj T\\\conj T-h_0\end{pmatrix}
=\mathbf v_1.
\]
这里用 $\dfrac{b(1-h_0^2)}{\kappa}=1+h_0^2$.
第二列分母在此点的导数是 $\dfrac{1+h_0^2}{2}$.
注意 $H$ 对非对角元的变号后, 其留数为
\[
\begin{aligned}
\Res_{z=\zeta_1}M_2
&=\dfrac{4\ii b(1-h_0^2)T}{(1+h_0^2)\Delta_-}
\begin{pmatrix}1-h_0\conj T\\\conj T-h_0\end{pmatrix}\\
&=\dfrac{4\ii\kappa T}{\Delta_-}
\begin{pmatrix}1-h_0\conj T\\\conj T-h_0\end{pmatrix}
=C_1\mathbf v_1.
\end{aligned}
\]

\paragraph{内部 $D_+$ 点 $\zeta_2$ 的正则列与留数.}
这里 $k=-\ii b$, 所以 $N=-2\ii bP$.
两列背景向量为 $(1,h_0^{-1})^T,(h_0^{-1},1)^T$.
所需投影乘法分别是
\[
\begin{aligned}
P\begin{pmatrix}1\\h_0^{-1}\end{pmatrix}
&=\dfrac{1-h_0^2}{h_0\Delta_-}
\begin{pmatrix}T-h_0\\1-h_0T\end{pmatrix},\\
P\begin{pmatrix}h_0^{-1}\\1\end{pmatrix}
&=\dfrac{(1-h_0^2)\conj T}{h_0\Delta_-}
\begin{pmatrix}T-h_0\\1-h_0T\end{pmatrix}.
\end{aligned}
\]
两式的标量来源可逐个查看:
\[
\begin{aligned}
(\conj T-h_0)+h_0^{-1}(1-h_0\conj T)
&=\dfrac{1-h_0^2}{h_0},\\
h_0^{-1}(\conj T-h_0)+(1-h_0\conj T)
&=\dfrac{(1-h_0^2)\conj T}{h_0}.
\end{aligned}
\]
正则第一列因此是
\[
M_1(\zeta_2)
=\dfrac{1+h_0^2}{h_0\Delta_-}
\begin{pmatrix}h_0-T\\1-h_0T\end{pmatrix}
=\mathbf v_2.
\]
第二列分母导数为 $\dfrac{1+h_0^2}{2h_0^2}$, 所以
\[
\begin{aligned}
\Res_{z=\zeta_2}M_2
&=\dfrac{4\ii b h_0(1-h_0^2)\conj T}
{(1+h_0^2)\Delta_-}
\begin{pmatrix}T-h_0\\-(1-h_0T)\end{pmatrix}\\
&=\dfrac{4\ii\kappa h_0\conj T}{\Delta_-}
\begin{pmatrix}T-h_0\\h_0T-1\end{pmatrix}
=C_2\mathbf v_2.
\end{aligned}
\]

\paragraph{外部 $D_-$ 点 $\conj{\zeta_1}$ 的正则列与留数.}
这里仍有 $N=-2\ii bP$, 但 $\lambda=-\ii\kappa$.
背景两列为 $(1,h_0)^T,(h_0,1)^T$.
直接相乘得到
\[
\begin{aligned}
P\begin{pmatrix}1\\h_0\end{pmatrix}
&=\dfrac{(1-h_0^2)\conj T}{\Delta_-}
\begin{pmatrix}T-h_0\\1-h_0T\end{pmatrix},\\
P\begin{pmatrix}h_0\\1\end{pmatrix}
&=\dfrac{1-h_0^2}{\Delta_-}
\begin{pmatrix}T-h_0\\1-h_0T\end{pmatrix}.
\end{aligned}
\]
正则第二列分母为 $-2\ii\kappa$, 因而
\[
M_2(\conj{\zeta_1})
=\dfrac{1+h_0^2}{\Delta_-}
\begin{pmatrix}h_0-T\\1-h_0T\end{pmatrix}
=\mathbf u_1.
\]
第一列分母导数为 $\dfrac{1+h_0^2}{2}$, 所以
\[
\begin{aligned}
\Res_{z=\conj{\zeta_1}}M_1
&=\dfrac{4\ii\kappa\conj T}{\Delta_-}
\begin{pmatrix}h_0-T\\1-h_0T\end{pmatrix}\\
&=-\conj{C_1}\mathbf u_1.
\end{aligned}
\]
其负号已经由 $\conj{C_1}=-\dfrac{4\ii\kappa\conj T}{1+h_0^2}$ 明确给出.

\paragraph{内部 $D_-$ 点 $\conj{\zeta_2}$ 的正则列与留数.}
此时 $N=2\ii b(I-P)$, $\lambda=-\ii\kappa$,
背景两列为 $(1,-h_0^{-1})^T,(-h_0^{-1},1)^T$.
直接计算
\[
\begin{aligned}
(I-P)\begin{pmatrix}1\\-h_0^{-1}\end{pmatrix}
&=\dfrac{(1-h_0^2)T}{h_0\Delta_-}
\begin{pmatrix}1-h_0\conj T\\-(\conj T-h_0)\end{pmatrix},\\
(I-P)\begin{pmatrix}-h_0^{-1}\\1\end{pmatrix}
&=\dfrac{1-h_0^2}{h_0\Delta_-}
\begin{pmatrix}-(1-h_0\conj T)\\\conj T-h_0\end{pmatrix}.
\end{aligned}
\]
正则第二列是
\[
M_2(\conj{\zeta_2})
=-\dfrac{1+h_0^2}{h_0\Delta_-}
\begin{pmatrix}1-h_0\conj T\\\conj T-h_0\end{pmatrix}
=\mathbf u_2.
\]
第一列分母导数为 $\dfrac{1+h_0^2}{2h_0^2}$, 给出
\[
\begin{aligned}
\Res_{z=\conj{\zeta_2}}M_1
&=\dfrac{4\ii\kappa h_0T}{\Delta_-}
\begin{pmatrix}1-h_0\conj T\\\conj T-h_0\end{pmatrix}\\
&=-\conj{C_2}\mathbf u_2.
\end{aligned}
\]
至此四个点全部按显式矩阵直接取过正则列和留数,
不再只依赖部分分式的抽象结论.

\paragraph{$N$ 的投影分解, 逆矩阵与行列式.}
利用 $I=P+(I-P)$, 原有 $N$ 可分解为
\[
N=(k-\ii b)P+(k+\ii b)(I-P).
\]
因为 $P(I-P)=(I-P)P=0$, 直接相乘得到
\[
\begin{aligned}
N^{-1}
&=\dfrac{P}{k-\ii b}+\dfrac{I-P}{k+\ii b}\\
&=\dfrac{(k+\ii b)P+(k-\ii b)(I-P)}{k^2+b^2}\\
&=\dfrac{(k-\ii b)I+2\ii bP}{k^2+b^2}.
\end{aligned}
\]
分母不为零处该式逐项成立, 其余处按亚纯函数理解.
两个互补秩一子空间上 $N$ 的特征值是 $k-\ii b,k+\ii b$,
所以
\[
\det N=k^2+b^2
=\lambda^2-q_0^2+b^2
=\lambda^2+\kappa^2.
\]
于是显式矩阵全部因子的行列式相乘为
\[
\begin{aligned}
\det M
&=\det H\,\det N\,\det Y
\dfrac{1}{(\lambda+\ii\kappa)(\lambda-\ii\kappa)}\det H^{-1}\\
&=\dfrac{(\lambda^2+\kappa^2)\gamma}{\lambda^2+\kappa^2}
=\gamma.
\end{aligned}
\]
在离散点, 该式是约去公因子后的亚纯恒等式.

\paragraph{两域约化为何由同一个有限系统固定.}
共轭变换对任意二阶矩阵的作用是
\[
\sigma_2\conj M\sigma_2
=\begin{pmatrix}
\conj{M_{22}}&-\conj{M_{21}}\\
-\conj{M_{12}}&\conj{M_{11}}
\end{pmatrix}.
\]
因此 $D_+$ 的第二列极点转成 $D_-$ 的第一列极点,
列交换时的负号正是负留数条件中的负号.
换叶关系的列形式是
\[
M_1(h(z))=\dfrac{\ii z}{q_0}M_2(z),\qquad
M_2(h(z))=\dfrac{\ii z}{q_0}M_1(z).
\]
例如 $h(\conj{\zeta_1})=\zeta_2$ 时,
$\dfrac{\ii\conj{\zeta_1}}{q_0}=h_0^{-1}$,
第一式正给出 $\mathbf v_2=\dfrac{\mathbf u_1}{h_0}$.
在 $h(\zeta_1)=\conj{\zeta_2}$ 时,
$\dfrac{\ii\zeta_1}{q_0}=-h_0^{-1}$,
第二式给出 $\mathbf u_2=-\dfrac{\mathbf v_1}{h_0}$.
这两个正则列关系已经由完整消元独立得到.
在留数中还要乘 $h'(z)=\dfrac{q_0^2}{z^2}$;
原有 $C$ 的配对公式恰好提供这个导数因子,
所以变换后的矩阵满足同一留数系统.
由于该系统可逆, 两域约化得以全局成立.

\paragraph{分支点与轮廓点的解析性.}
在 $z=\pm\ii q_0$ 处 $\lambda=0$,
$\lambda\pm\ii\kappa$ 都非零, $\det N=\kappa^2>0$.
只有 $Y$ 的行列式为零, 且
\[
Y(\ii q_0)=I-\sigma_1,\qquad
Y(-\ii q_0)=I+\sigma_1.
\]
这两个矩阵均非零且秩为一, 所以 $M$ 解析且秩为一.
在交点 $z=\pm q_0$ 处, $\lambda=\pm q_0$ 为非零实数,
列分母也不为零.
在共振原像 $z_a^\pm$ 处, $\lambda=\pm\sqrt{a^2+q_0^2}$,
列分母非零且 $\gamma>0$, 所以 $M$ 解析并可逆.
全部判断直接来自显式有理式的分母, 无需附加轮廓上的跳跃论证.
% END EXPANSION 08-H
重构公式将同一矩阵的无穷远系数与共振值分别转化为光场和介质. 以下结论同时给出三个物理场及其 Lax 相容性.

\begin{theorem}[单轨道三场解]\label{thm:direct-HMB}\label{prop:oneboundary}
设定理 \ref{thm:oneRHP} 的条件成立, $N$ 按 \eqref{eq:Nfactor} 定义, $n_{jk}=N(a)_{jk}$. 则物理场
\begin{equation}\label{eq:finalfields}
\begin{aligned}
\widetilde q&=-e^{-\ii\Omega t}(q_0+4bP_{12}),\\
\widetilde p&=\dfrac{e^{-\ii\Omega t}(N(a)R_\infty N(a)^\dagger)_{12}}{a^2+b^2},\qquad
\eta=\dfrac{(N(a)R_\infty N(a)^\dagger)_{11}}{a^2+b^2}
\end{aligned}
\end{equation}
在全体实 $x,t$ 上光滑, 满足原 HMB 系统 \eqref{eq:HMB}、Bloch 约束 \eqref{eq:C} 及两端同相边界 \eqref{eq:targetboundary}. 两个共振原像重构出同一个 Hermitian 迹零介质矩阵. 介质分量为
\begin{align}
\eta&=\dfrac{\eta_0}{a^2+b^2}\left[|n_{11}|^2-|n_{12}|^2
+\dfrac{\ii q_0}{a}(n_{11}\conj{n_{12}}-n_{12}\conj{n_{11}})\right],\label{eq:etacomponent}\\
e^{\ii\Omega t}\widetilde p&=\dfrac{\eta_0}{a^2+b^2}\left[n_{11}\conj{n_{21}}-n_{12}\conj{n_{22}}
+\dfrac{\ii q_0}{a}(n_{11}\conj{n_{22}}-n_{12}\conj{n_{21}})\right].\label{eq:pcomponent}
\end{align}
固定 $t$ 时, 三个场的端态误差为 $O(e^{-2\kappa|x|})$.
\end{theorem}
\begin{proof}
在无穷远展开 \eqref{eq:explicitM}, 得
\[
M^{(1)}=H(2\ii bI-4\ii bP-\ii q_0\sigma_1-2\ii\kappa\sg)H^{-1}.
\]
于是规范光场 $q=\ii M^{(1)}_{12}=-q_0-4bP_{12}$. 在实共振点, 由 $Y\lambda\sg Y^{-1}=\ii U_\infty$ 及 $N(a)^{-1}=\dfrac{N(a)^\dagger}{a^2+b^2}$,
\begin{equation}\label{eq:direct-medium}
M(z_a)\dfrac{\eta_0\lambda(z_a)}{a}\sg M(z_a)^{-1}
=\dfrac{HN(a)R_\infty N(a)^\dagger H^{-1}}{a^2+b^2}=:R.
\end{equation}
右端与谱叶无关. $N(a)N(a)^\dagger=(a^2+b^2)I$, 因而该酉共轭保持 Hermitian 性、迹零及 $R_\infty$ 的特征值. 特别地,
$R^2=\eta_0^2\left(1+\dfrac{q_0^2}{a^2}\right)I$, 给出 Bloch 约束. 结合 \eqref{eq:original-fields} 得到 \eqref{eq:finalfields}; 矩阵相乘给出两条分量式.

当 $x\to\pm\infty$ 时, $|T|$ 分别趋于无穷和零, 因而
\begin{equation}\label{eq:Pends}
P_+=\dfrac{1}{1+h_0^2}\begin{pmatrix}1&-h_0\\-h_0&h_0^2\end{pmatrix},\qquad
P_-=\dfrac{1}{1+h_0^2}\begin{pmatrix}h_0^2&-h_0\\-h_0&1\end{pmatrix}.
\end{equation}
这给出 $q\to q_0$. 将端态代入 $N$, 有
\[
N_+Y=Y_{-q_0}\diag(\lambda-\ii\kappa,\lambda+\ii\kappa),\qquad
N_-Y=Y_{-q_0}\diag(\lambda+\ii\kappa,\lambda-\ii\kappa),
\]
其中 $Y_{-q_0}=I+\dfrac{\ii q_0}{z}\sigma_1$. 利用 $HY_{-q_0}H^{-1}=Y$, 得
\begin{equation}\label{eq:direct-M-ends}
M\longrightarrow Y\quad(x\to-\infty),\qquad
M\longrightarrow Y\diag(s_{22},s_{22}^{-1})\quad(x\to+\infty).
\end{equation}
对角因子与 $\sg$ 对易, 所以两端 $R\to R_\infty$. 由 $|T|=e^{2\kappa[x-\operatorname{Re}F(\ii b)t]}$, 端态误差及任意固定阶导数均按 $e^{-2\kappa|x|}$ 衰减, 对紧 $t$ 区间一致. $\Delta_->0$ 与 $a^2+b^2>0$ 保证全局光滑.

对避开特殊点的 $z$ 定义
\begin{equation}\label{eq:inverse-log-lax}
U=(M_x-\ii\lambda M\sg)M^{-1},\qquad
V=(M_t+\ii\lambda F M\sg)M^{-1}.
\end{equation}
它们是 $\phi=Me^{\ii\theta\sg}$ 的两个对数导数. 换叶时 $\lambda$ 和 $\theta$ 变号, $\sigma_1$ 与 $\sg$ 反对易, 故两矩阵在 $z$ 和 $h(z)$ 取值相同, 是 $k$ 的单值有理矩阵.

在 $\zeta_j$ 附近, $\phi$ 的列留数为 $\Res\phi_2=c_j\phi_1(\zeta_j)$. 常数 $c_j$ 与 $x,t$ 无关, 故
\[
\phi=H_j(z)\left(I+\dfrac{c_je_1e_2^T}{z-\zeta_j}\right),\qquad
\det H_j(\zeta_j)=\gamma(\zeta_j)\ne0.
\]
右因子不依赖 $x,t$, 因而两个对数导数的离散极点可去; 共轭极点同理.

在分支点 $z_*=\pm\ii q_0$, 换叶约化使 $M(z_*)$ 消灭 $C(z_*)$ 的固定 $-1$ 特征空间. 该核不依赖 $x,t$, 因此 $M_x(z_*)$ 和 $M_t(z_*)$ 也消灭它. 又 $\lambda(z_*)=0$, $F$ 在此有限, \eqref{eq:inverse-log-lax} 的两个分子均消灭同一核. $\det M=\gamma$ 在 $z_*$ 有一阶零点, 而 $\operatorname{adj}M(z_*)$ 的像正是该核, 故两个商的分子抵消行列式的零点. 这证明分支点可去. 轮廓交点处 $\gamma\ne0$, 没有逆矩阵奇性. 在共振点, $U$ 解析, $V$ 的唯一留数为
\[
\Res_{k=a}V=\dfrac{\ii\eta_0\lambda(z_a)}{a}M(z_a)\sg M(z_a)^{-1}=\ii R.
\]

在第一叶无穷远, $M=I+O(k^{-1})$, $M_x,M_t=O(k^{-1})$. 由 \eqref{eq:reconq},
\[
U=-\ii k\sg+Q+O(k^{-1}).
\]
有限点均可去, Liouville 定理给出 $U=-\ii k\sg+Q$. 恒等式
\[
V=-FU+(M_t+FM_x)M^{-1}
\]
的第二项至多一次增长, 因而
\begin{equation}\label{eq:inverse-V-form}
V=4\ii\beta k^3\sg+k^2(-\ii\alpha\sg-4\beta Q)
+kW_1+W_0+\dfrac{\ii R}{k-a}.
\end{equation}
由 $\det\phi=\gamma$ 与 $x,t$ 无关, 两矩阵均迹零. 实谱上的共轭约化给出 $M^\dagger M=\gamma I$, 因此 $U,V$ 在实 $k\ne a$ 上反 Hermitian, 各系数具有相应约化.

作为同一 $\phi$ 的对数导数, 两矩阵恒满足零曲率关系. 比较 $k^2$ 项的非对角元及 $k$ 项的对角元, 得
\[
\begin{aligned}
(W_1)_{12}&=\alpha q-2\ii\beta q_x,\\
\partial_x(W_1)_{11}&=-2\ii\beta\partial_x|q|^2.
\end{aligned}
\]
故 $(W_1)_{11}=-2\ii\beta|q|^2+\chi_1(t)$. 由 \eqref{eq:direct-M-ends}, 两端 $V\to-FU_\infty$, 从而 $\chi_1(t)=0$. 再比较 $k$ 项的非对角元及常数项的对角元, 有
\[
\begin{aligned}
(W_0)_{12}&=2\beta|q|^2q+\dfrac{\ii\alpha}{2}q_x+\beta q_{xx},\\
\partial_x(W_0)_{11}&=\partial_x\left[\dfrac{\ii\alpha}{2}|q|^2-\beta(q\conj{q_x}-q_x\conj{q})\right].
\end{aligned}
\]
同一端态固定后一积分常数为 $\dfrac{\ii\Omega}{2}$. 于是 $W_1=V_1(q)$、$W_0=V_0(q)+\dfrac{\ii\Omega}{2}\sg$, 其余分量由迹零和约化确定, 得到完整规范 Lax 对 \eqref{eq:gaugelax}.

零曲率的共振留数和常数非对角元分别给出
\[
R_x=[-\ii a\sg+Q,R],\qquad
q_t=\ii\alpha\left(\dfrac{q_{xx}}{2}+|q|^2q\right)
+\beta(q_{xxx}+6|q|^2q_x)+2p+\ii\Omega q.
\]
前式为两条介质方程; 恢复 \eqref{eq:original-fields} 即得原 HMB 系统.

最后, 在普通连续谱上令
\[
\phi_-=Me^{\ii\theta\sg},\qquad
\phi_+=M\diag(s_{22}^{-1},s_{22})e^{\ii\theta\sg}.
\]
已恢复的 Lax 对及 \eqref{eq:direct-M-ends} 表明两者具有相应端点归一化, 由 Jost 唯一性即为重构场的 Jost 解. 所以 $S=\diag(s_{22}^{-1},s_{22})$, 反射为零, 透射函数为 \eqref{eq:onezeros}; 外部零点的列关系给出 $b_1=c_1s_{22}'(\zeta_1)=1$.
\end{proof}

% BEGIN EXPANSION 08-I
\subsubsection*{三个物理场的重构乘法, 端态矩阵与误差}
无穷远处, 三个谱因子的展开分别是
\[
\begin{aligned}
N(k(z))
&=\dfrac{z}{2}I+\ii bI-2\ii bP
-\dfrac{q_0^2}{2z}I,\\
Y(z)&=I-\dfrac{\ii q_0}{z}\sigma_1,\\
\diag\left(\dfrac{1}{\lambda+\ii\kappa},
\dfrac{1}{\lambda-\ii\kappa}\right)
&=\dfrac{2}{z}
\left(I-\dfrac{2\ii\kappa}{z}\sg+O(z^{-2})\right).
\end{aligned}
\]
先乘前两个矩阵, 再乘第三个, 得
\[
\begin{aligned}
NY
&=\dfrac{z}{2}I
+\ii bI-2\ii bP-\dfrac{\ii q_0}{2}\sigma_1
+O(z^{-1}),\\
NY\diag\left(\dfrac{1}{\lambda+\ii\kappa},
\dfrac{1}{\lambda-\ii\kappa}\right)
&=I+\dfrac{2\ii bI-4\ii bP-\ii q_0\sigma_1-2\ii\kappa\sg}{z}
+O(z^{-2}).
\end{aligned}
\]
最后乘 $H,H^{-1}$, 非对角元变号, 所以
\[
M^{(1)}_{12}=4\ii bP_{12}+\ii q_0,
\qquad
\ii M^{(1)}_{12}=-4bP_{12}-q_0.
\]
这也可由有限求和独立核验:
\[
\begin{aligned}
C_1(\mathbf v_1)_1+C_2(\mathbf v_2)_1
&=\dfrac{4\ii\kappa}{\Delta_-}
\bigl(T(1-h_0\conj T)+h_0\conj T(T-h_0)\bigr)\\
&=\dfrac{4\ii\kappa}{\Delta_-}(T-h_0^2\conj T),\\
q
&=q_0-\dfrac{4\kappa(T-h_0^2\conj T)}{\Delta_-}.
\end{aligned}
\]
上式与 $-q_0-4bP_{12}$ 相等, 因为通分后需要核验的恒等式只是
\[
\begin{aligned}
2q_0\Delta_-
={}&4\kappa(T-h_0^2\conj T)
-4b(T-h_0)(1-h_0\conj T).
\end{aligned}
\]
代入 $2b=\dfrac{q_0(1+h_0^2)}{h_0}$ 与
$2\kappa=\dfrac{q_0(1-h_0^2)}{h_0}$ 后, 右侧展开为
\[
2q_0(1+h_0^2)(1+|T|^2)-4q_0h_0(T+\conj T),
\]
正是左侧.

\paragraph{共振矩阵的消因子与酉共轭.}
在 $k=a$ 上, $P^\dagger=P$ 给出
\[
N(a)^\dagger=(a-\ii b)I+2\ii bP.
\]
直接乘法为
\[
\begin{aligned}
N(a)N(a)^\dagger
={}&(a+\ii b)(a-\ii b)I\\
&+2\ii b(a+\ii b)P
-2\ii b(a-\ii b)P
+4b^2P^2\\
={}&(a^2+b^2)I-4b^2P+4b^2P\\
={}&(a^2+b^2)I.
\end{aligned}
\]
所以 $N(a)^{-1}=\dfrac{N(a)^\dagger}{a^2+b^2}$.
共振重构中, 右侧的列分母矩阵与 $\sg$ 都是对角矩阵,
故该分母矩阵及其逆矩阵在共轭乘积中相消.
同时 $H^{-1}\sg H=\sg$.
于是
\[
\begin{aligned}
M\dfrac{\eta_0\lambda}{a}\sg M^{-1}
&=HN Y\dfrac{\eta_0\lambda}{a}\sg Y^{-1}N^{-1}H^{-1}\\
&=HN\dfrac{\ii\eta_0}{a}U_\infty(a)N^{-1}H^{-1}\\
&=HN R_\infty N^{-1}H^{-1}\\
&=\dfrac{HN R_\infty N^\dagger H^{-1}}{a^2+b^2}.
\end{aligned}
\]
第一步消去了所有谱叶依赖的列分母;
第二步使用 $Y\lambda\sg Y^{-1}=\ii U_\infty$.
所得右侧仅含 $a$, 所以两共振原像自动给出同一矩阵.

\paragraph{介质分量中每一个乘积的来源.}
原有 $n_{jk}=N(a)_{jk}$, 因而
\[
N(a)=\begin{pmatrix}n_{11}&n_{12}\\n_{21}&n_{22}\end{pmatrix},\qquad
N(a)^\dagger=
\begin{pmatrix}\conj{n_{11}}&\conj{n_{21}}\\
\conj{n_{12}}&\conj{n_{22}}\end{pmatrix}.
\]
第一轮矩阵乘法为
\[
N(a)R_\infty
=\eta_0\begin{pmatrix}
n_{11}-\dfrac{\ii q_0}{a}n_{12}
&\dfrac{\ii q_0}{a}n_{11}-n_{12}\\
n_{21}-\dfrac{\ii q_0}{a}n_{22}
&\dfrac{\ii q_0}{a}n_{21}-n_{22}
\end{pmatrix}.
\]
第二轮右乘 $N(a)^\dagger$ 的四元分别为
\[
\begin{aligned}
(N R_\infty N^\dagger)_{11}
={}&\eta_0\left[
\left(n_{11}-\dfrac{\ii q_0}{a}n_{12}\right)\conj{n_{11}}
+\left(\dfrac{\ii q_0}{a}n_{11}-n_{12}\right)\conj{n_{12}}
\right]\\
={}&\eta_0\left[
|n_{11}|^2-|n_{12}|^2
+\dfrac{\ii q_0}{a}
(n_{11}\conj{n_{12}}-n_{12}\conj{n_{11}})
\right],\\
(N R_\infty N^\dagger)_{12}
={}&\eta_0\left[
\left(n_{11}-\dfrac{\ii q_0}{a}n_{12}\right)\conj{n_{21}}
+\left(\dfrac{\ii q_0}{a}n_{11}-n_{12}\right)\conj{n_{22}}
\right]\\
={}&\eta_0\left[
n_{11}\conj{n_{21}}-n_{12}\conj{n_{22}}
+\dfrac{\ii q_0}{a}
(n_{11}\conj{n_{22}}-n_{12}\conj{n_{21}})
\right],\\
(N R_\infty N^\dagger)_{21}
={}&\eta_0\left[
\left(n_{21}-\dfrac{\ii q_0}{a}n_{22}\right)\conj{n_{11}}
+\left(\dfrac{\ii q_0}{a}n_{21}-n_{22}\right)\conj{n_{12}}
\right],\\
(N R_\infty N^\dagger)_{22}
={}&\eta_0\left[
\left(n_{21}-\dfrac{\ii q_0}{a}n_{22}\right)\conj{n_{21}}
+\left(\dfrac{\ii q_0}{a}n_{21}-n_{22}\right)\conj{n_{22}}
\right].
\end{aligned}
\]
$H$ 的共轭作用将此矩阵的 $(1,2)$ 元变号,
而从 $R_{12}$ 恢复 $\widetilde p$ 时又有一个负号,
这两个负号抵消, 给出 \eqref{eq:pcomponent}.
对角元不变, 直接给出 \eqref{eq:etacomponent}.
介质矩阵的 Hermitian 性由
\[
(HN R_\infty N^\dagger H^{-1})^\dagger
=HN R_\infty N^\dagger H^{-1}
\]
得到, 因为 $H^{-1}=H^\dagger$ 且 $R_\infty^\dagger=R_\infty$.
迹用循环性与 $N^\dagger N=(a^2+b^2)I$ 得
\[
\tr R
=\dfrac{\tr(R_\infty N^\dagger N)}{a^2+b^2}
=\tr R_\infty=0.
\]
背景矩阵平方的两个对角元均为
$\eta_0^2+\dfrac{\eta_0^2q_0^2}{a^2}$, 两个非对角元为零:
\[
R_\infty^2
=\eta_0^2\left(1+\dfrac{q_0^2}{a^2}\right)I.
\]
把 $N^\dagger N$ 在 $R^2$ 的中间约掉, 就得到同一平方恒等式.
又因
\[
R=\begin{pmatrix}\eta&-p\\-\conj p&-\eta\end{pmatrix},\qquad
R^2=(\eta^2+|p|^2)I,
\]
且 $|p|=|\widetilde p|$, 所以得到 Bloch 约束.

\paragraph{两端的 $N$ 与背景列的逐元相乘.}
由 $\Delta_-$ 的最高次项与常数项分别读取 $P_+,P_-$.
再用 $2bh_0=q_0(1+h_0^2)$, 得
\[
N_+=\begin{pmatrix}k-\ii\kappa&\ii q_0\\
\ii q_0&k+\ii\kappa\end{pmatrix},\qquad
N_-=\begin{pmatrix}k+\ii\kappa&\ii q_0\\
\ii q_0&k-\ii\kappa\end{pmatrix}.
\]
例如 $N_+Y$ 的四个元是
\[
\begin{aligned}
(N_+Y)_{11}&=k-\ii\kappa+\dfrac{q_0^2}{z}=\lambda-\ii\kappa,\\
(N_+Y)_{12}&=\ii q_0-\dfrac{\ii q_0}{z}(k-\ii\kappa)
=\dfrac{\ii q_0}{z}(\lambda+\ii\kappa),\\
(N_+Y)_{21}&=\ii q_0-\dfrac{\ii q_0}{z}(k+\ii\kappa)
=\dfrac{\ii q_0}{z}(\lambda-\ii\kappa),\\
(N_+Y)_{22}&=k+\ii\kappa+\dfrac{q_0^2}{z}=\lambda+\ii\kappa.
\end{aligned}
\]
即 $Y_{-q_0}\diag(\lambda-\ii\kappa,\lambda+\ii\kappa)$.
交换 $\kappa$ 的正负, 就逐元得到
\[
N_-Y=Y_{-q_0}\diag(\lambda+\ii\kappa,\lambda-\ii\kappa).
\]
乘回显式矩阵中的列分母, 再用 $HY_{-q_0}H^{-1}=Y$,
正是 \eqref{eq:direct-M-ends}.
因为两端的对角因子与 $\sg$ 对易, 它们在介质共轭中完全相消.
因此
\[
R\longrightarrow Y\dfrac{\eta_0\lambda}{a}\sg Y^{-1}=R_\infty.
\]
光场端态也可由投影直接读出:
\[
q\longrightarrow -q_0-4b\left(-\dfrac{h_0}{1+h_0^2}\right)
=-q_0+2q_0=q_0.
\]

\paragraph{指数误差怎样从有理分子分母中读出.}
在 $x\to-\infty$ 时 $|T|\to0$,
$\Delta_-=(1+h_0^2)+O(|T|)$,
$P$ 的各元分子是其常数项加 $O(|T|)$.
因此
\[
P-P_-=O(|T|).
\]
在 $x\to+\infty$ 时, 将分子分母同时除以 $|T|^2$,
分母成为 $(1+h_0^2)+O(|T|^{-1})$,
各分子成为端态常数加 $O(|T|^{-1})$, 所以
\[
P-P_+=O(|T|^{-1}).
\]
而
\[
|T|=e^{2\kappa[x-\operatorname{Re}F(\ii b)t]},
\]
故固定 $t$ 后两端分别是 $O(e^{2\kappa x})$ 与
$O(e^{-2\kappa x})$, 合写为 $O(e^{-2\kappa|x|})$.
因为
\[
T_x=2\kappa T,\quad
T_t=-2\kappa F(\ii b)T,\quad
(\conj T)_x=2\kappa\conj T,\quad
(\conj T)_t=-2\kappa\conj{F(\ii b)}\conj T,
\]
任何固定阶 $x,t$ 导数仍具有相同端态衰减阶.
分母始终受 $\Delta_-\ge(1-h_0)^2(1+|T|^2)>0$ 控制,
所以求导不产生真实奇点.
$q$ 是 $P_{12}$ 的仿射函数,
$R$ 是 $P$ 的多项式除以正常数 $a^2+b^2$,
从而三场及任意固定阶导数都得到同一误差阶.
% END EXPANSION 08-I
写 $s=\xi+\ii\nu$, 光场公式中的投影元还可表示为
\begin{equation}\label{eq:hyperbolic-projector}
P_{12}=\dfrac{-2h_0\cosh(2\xi)+(1+h_0^2)\cos(2\nu)+\ii(1-h_0^2)\sin(2\nu)}
{2(1+h_0^2)\cosh(2\xi)-4h_0\cos(2\nu)}.
\end{equation}
分母不小于 $2(1-h_0)^2$, 局域包络与内部振荡分别由 $\xi$ 和 $\nu$ 控制.
% BEGIN EXPANSION 08-J
\subsubsection*{对数导数的每种潜在极点与谱变量单值性}
先直接微分 $\phi=Me^{\ii\theta\sg}$.
因为 $\theta_x=-\lambda$ 与 $\theta_t=\lambda F$,
\[
\begin{aligned}
\phi_x
&=(M_x-\ii\lambda M\sg)e^{\ii\theta\sg},\\
\phi_t
&=(M_t+\ii\lambda F M\sg)e^{\ii\theta\sg}.
\end{aligned}
\]
右乘 $\phi^{-1}=e^{-\ii\theta\sg}M^{-1}$,
就分别得到 $ U, V$.
以下要检查 $M^{-1}$ 产生的潜在奇性, 不能只检查 $M$ 的分母.

\paragraph{离散点处, 不依赖 $x,t$ 的右因子怎样消去极点.}
在 $\zeta_j$ 处, $\phi_1$ 正则,
且 $\phi_2$ 的留数为 $c_j\phi_1(\zeta_j)$.
利用 $(e_1e_2^T)^2=e_1(e_2^Te_1)e_2^T=0$,
原证明的正则因子可以直接定义为
\[
 H_j(z)
=\phi(z)\left(I-\dfrac{c_je_1e_2^T}{z-\zeta_j}\right).
\]
其两列是
\[
( H_j)_1=\phi_1(z),\qquad
( H_j)_2=\phi_2(z)-\dfrac{c_j\phi_1(z)}{z-\zeta_j}.
\]
将 $\phi_1(z)=\phi_1(\zeta_j)+O(z-\zeta_j)$ 代入第二列,
其负一次项恰好相消, 所以 $ H_j$ 解析.
右因子的行列式为 1, 因而
\[
\det H_j(z)=\det\phi(z)=\gamma(z).
\]
由于离散点不在分支点, $\gamma(\zeta_j)\ne0$,
故 $ H_j$ 在该点可逆.
再将反因子乘回去, 有
\[
\phi= H_j
\left(I+\dfrac{c_je_1e_2^T}{z-\zeta_j}\right).
\]
因为 $c_j,\zeta_j$ 与 $x,t$ 无关, 对 $x$ 微分时只微分 $ H_j$:
\[
\begin{aligned}
\phi_x\phi^{-1}
&=( H_j)_x
\left(I+\dfrac{c_je_1e_2^T}{z-\zeta_j}\right)
\left(I-\dfrac{c_je_1e_2^T}{z-\zeta_j}\right)
 H_j^{-1}\\
&=( H_j)_x H_j^{-1}.
\end{aligned}
\]
对 $t$ 微分得到同样的式子.
所以 $ U, V$ 的离散极点都可去.
在共轭极点处, 用 $-\conj{c_j}e_2e_1^T$ 替换
$c_je_1e_2^T$, 同样有平方为零的恒等式,
因此两组离散点全部消去.

\paragraph{两个分支点处固定核的具体向量.}
在 $z_* =\ii q_0$ 处, 换叶标量矩阵是
\[
\dfrac{\ii z_*}{q_0}\sigma_1=-\sigma_1.
\]
$h(z_*)=z_*$, 所以 $M(z_*)=M(z_*)(-\sigma_1)$,
即 $M(z_*)(I+\sigma_1)=0$.
因此
\[
M(\ii q_0)\begin{pmatrix}1\\1\end{pmatrix}=0.
\]
在 $z_*=-\ii q_0$ 处, 同理得到
\[
M(-\ii q_0)\begin{pmatrix}1\\-1\end{pmatrix}=0.
\]
这两条核向量是常向量, 不依赖 $x,t$.
对相应等式求 $x,t$ 导数, 就得到
\[
\begin{aligned}
M_x(\ii q_0)\begin{pmatrix}1\\1\end{pmatrix}
&=M_t(\ii q_0)\begin{pmatrix}1\\1\end{pmatrix}=0,\\
M_x(-\ii q_0)\begin{pmatrix}1\\-1\end{pmatrix}
&=M_t(-\ii q_0)\begin{pmatrix}1\\-1\end{pmatrix}=0.
\end{aligned}
\]
此处 $M$ 的秩为一, 所以其伴随矩阵也是秩一且非零.
由 $M\operatorname{adj}M=(\det M)I=0$,
$\operatorname{adj}M(z_*)$ 的每一列都属于上述一维核.
又因 $\lambda(z_*)=0$ 且 $F(k(z_*))$ 有限,
\[
\begin{aligned}
\bigl(M_x-\ii\lambda M\sg\bigr)(z_*)
\operatorname{adj}M(z_*)&=0,\\
\bigl(M_t+\ii\lambda F M\sg\bigr)(z_*)
\operatorname{adj}M(z_*)&=0.
\end{aligned}
\]
而
\[
\gamma'(z_*)=-\dfrac{2q_0^2}{z_*^3}\ne0.
\]
所以将 $M^{-1}=\dfrac{\operatorname{adj}M}{\gamma}$ 代入对数导数后,
分子与分母都至少含一因子 $z-z_*$, 且分母恰为一阶零.
约去该因子就证明两个分支点均可去.
仅知道 $\det M=0$ 并不足够;
关键是上面已经逐一确认导数分子也消灭同一个固定核.

\paragraph{共振点与轮廓交点.}
轮廓交点处 $M$ 解析且 $\gamma\ne0$,
所以其逆矩阵也解析, 两对数导数均无奇性.
在 $k=a$ 处, $M,M_x,M_t$ 都解析且 $M$ 可逆,
$ U$ 因而解析.
在 $ V$ 中仅 $F$ 的共振项产生极点:
\[
F(k)=4\beta k^2-\alpha k-2\beta q_0^2
+\dfrac{\eta_0}{a(k-a)}.
\]
因此按 $k$ 取留数得到
\[
\begin{aligned}
\Res_{k=a} V
&=\lim\limits_{k\to a}(k-a)
\ii\lambda F M\sg M^{-1}\\
&=\dfrac{\ii\eta_0\lambda(z_a)}{a}
M(z_a)\sg M(z_a)^{-1}
=\ii R.
\end{aligned}
\]
这里取的是 $k$ 留数, 没有混入 $z$ 留数的导数因子.
两个 $z$ 原像因换叶重构一致而给出同一个 $k$ 留数.

\paragraph{换叶不变怎样消除 $k$ 平面的多值性.}
换叶时
\[
M(h(z))=M(z)\dfrac{\ii z}{q_0}\sigma_1,
\quad \lambda(h(z))=-\lambda(z),
\quad k(h(z))=k(z).
\]
其中右乘因子与 $x,t$ 无关, 所以
\[
\begin{aligned}
M_x(h(z))M(h(z))^{-1}&=M_x(z)M(z)^{-1},\\
M_t(h(z))M(h(z))^{-1}&=M_t(z)M(z)^{-1}.
\end{aligned}
\]
带 $\lambda$ 的项中, 又有
\[
\begin{aligned}
\lambda(h(z))M(h(z))\sg M(h(z))^{-1}
&=-\lambda(z)M(z)\sigma_1\sg\sigma_1M(z)^{-1}\\
&=\lambda(z)M(z)\sg M(z)^{-1}.
\end{aligned}
\]
$F$ 仅依赖 $k$, 故换叶不变.
由此逐项得到
\[
 U(h(z))= U(z),\qquad
 V(h(z))= V(z).
\]
显式 $M$ 对 $z$ 有理, 对 $x,t$ 微分只改变其系数,
因此这两个矩阵对 $z$ 都是有理矩阵.
关系 $z^2-2kz-q_0^2=0$ 说明 $z$ 是关于 $k$ 的二次代数变量.
任一关于 $z$ 的有理函数均可在该关系下化为
关于 $k$ 的有理函数加上 $z$ 乘以关于 $k$ 的有理函数.
再比较其在 $z$ 与 $h(z)$ 的值,
差值等于第二个系数乘以
\[
z-h(z)=z+\dfrac{q_0^2}{z}=2\lambda.
\]
换叶不变迫使第二个系数恒为零.
因此两矩阵实际上都是 $k$ 的单值有理矩阵.
这里没有把在分支点消去 $z$ 奇性误当成自动消去多值性;
两件事分别由固定核与换叶不变证明.

\paragraph{无穷远次数与空间 Lax 矩阵的恢复.}
在第一叶 $k\to\infty$ 时,
\[
\lambda=k+\dfrac{q_0^2}{2k}+O(k^{-3}),\qquad
M=I+O(k^{-1}),\qquad M_x,M_t=O(k^{-1}).
\]
为把常数项也逐元读出, 保留原来的 $z^{-1}$ 系数:
\[
M=I+\dfrac{M^{(1)}}{z}+O(z^{-2}),\qquad
M^{-1}=I-\dfrac{M^{(1)}}{z}+O(z^{-2}).
\]
于是按矩阵先后顺序相乘,
\begin{align*}
M\sg M^{-1}
&=\sg+\dfrac{M^{(1)}\sg-\sg M^{(1)}}{z}+O(z^{-2})\\
&=\sg-\dfrac{[\sg,M^{(1)}]}{z}+O(z^{-2}),\\
(M_x-\ii\lambda M\sg)M^{-1}
&=-\ii\lambda\sg+\dfrac{\ii\lambda}{z}[\sg,M^{(1)}]+O(z^{-1})\\
&=-\ii k\sg+\dfrac{\ii}{2}[\sg,M^{(1)}]+O(k^{-1}).
\end{align*}
最后一行使用 $\lambda-k=O(k^{-1})$ 及 $\dfrac{\lambda}{z}=\dfrac{1}{2}+O(z^{-2})$.
这里交换子的四个元是
\[
[\sg,M^{(1)}]
=\begin{pmatrix}0&2M^{(1)}_{12}\\-2M^{(1)}_{21}&0\end{pmatrix},\qquad
\dfrac{\ii}{2}[\sg,M^{(1)}]
=\begin{pmatrix}0&\ii M^{(1)}_{12}\\-\ii M^{(1)}_{21}&0\end{pmatrix}.
\]
由 $q=\ii M^{(1)}_{12}$ 以及共轭约化
$M^{(1)}_{21}=-\conj{M^{(1)}_{12}}$, 右上元就是 $q$,
左下元就是 $-\conj q$, 故常数矩阵正是 $Q(q)$.
于是
\[
 U=-\ii k\sg+Q+O(k^{-1}).
\]
已经证明 $ U$ 的全部有限点奇性可去.
故每个矩阵元减去 $(-\ii k\sg+Q)$ 的对应元后,
是无穷远趋零的整函数, 必恒为零:
\[
 U=-\ii k\sg+Q.
\]
对于时间矩阵, 精确重排是
\[
\begin{aligned}
-F U+(M_t+FM_x)M^{-1}
&=-F(M_x-\ii\lambda M\sg)M^{-1}
+(M_t+FM_x)M^{-1}\\
&=(M_t+\ii\lambda F M\sg)M^{-1}= V.
\end{aligned}
\]
$F=O(k^2)$ 与 $M_x=O(k^{-1})$ 使最后的修正项至多为 $O(k)$.
因此三次项和二次项完全由 $-F U$ 确定:
\[
\begin{aligned}
-F(-\ii k\sg+Q)
={}&4\ii\beta k^3\sg
+k^2(-\ii\alpha\sg-4\beta Q)+O(k).
\end{aligned}
\]
减去这些项和 $\dfrac{\ii R}{k-a}$ 后, 仅余一次多项式,
于是得到原文的 $kW_1+W_0$.

\paragraph{迹零与反 Hermitian 约化的独立元.}
指数矩阵的行列式是 1, 所以 $\det\phi=\gamma$.
对行列式微分的公式给出
\[
\tr(\phi_x\phi^{-1})=\partial_x\log\gamma=0,\qquad
\tr(\phi_t\phi^{-1})=\partial_t\log\gamma=0.
\]
在实 $z$ 上, 共轭约化逐元为
\[
M_{22}=\conj{M_{11}},\qquad
M_{21}=-\conj{M_{12}}.
\]
因此
\[
\begin{aligned}
M^\dagger M
&=\begin{pmatrix}
\conj{M_{11}}&-M_{12}\\
\conj{M_{12}}&M_{11}
\end{pmatrix}
\begin{pmatrix}
M_{11}&M_{12}\\
-\conj{M_{12}}&\conj{M_{11}}
\end{pmatrix}\\
&=(|M_{11}|^2+|M_{12}|^2)I=\gamma I.
\end{aligned}
\]
又 $\theta$ 在实谱上为实数, 故 $\phi^\dagger\phi=\gamma I$.
对 $x$ 或 $t$ 微分, 用 $\gamma_x=\gamma_t=0$,
便得到 $ U^\dagger=- U$ 与
$ V^\dagger=- V$.
因此每个多项式系数矩阵都由纯虚的 $(1,1)$ 元及复数 $(1,2)$ 元决定:
\[
(W_1)_{22}=-(W_1)_{11},\quad
(W_1)_{21}=-\conj{(W_1)_{12}},
\]
\[
(W_0)_{22}=-(W_0)_{11},\quad
(W_0)_{21}=-\conj{(W_0)_{12}}.
\]

\subsubsection*{零曲率逐次幂恢复全部 Lax 系数与积分常数}
由同一个可逆 $\phi$ 定义两个对数导数,
混合偏导交换给出
\[
\begin{aligned}
 U_t
&=\phi_{xt}\phi^{-1}- U V,\\
 V_x
&=\phi_{tx}\phi^{-1}- V U.
\end{aligned}
\]
相减便是
\[
 U_t- V_x+[ U, V]=0.
\]
下面逐个比较矩阵元, 不以已经知道 HMB 方程作为前提.

\paragraph{$k^2$ 项的非对角元.}
这一项的矩阵等式是
\[
4\beta Q_x+[-\ii\sg,W_1]+[Q,-\ii\alpha\sg-4\beta Q]=0.
\]
其中
\[
[Q,-4\beta Q]=0,\qquad
[Q,-\ii\alpha\sg]_{12}=2\ii\alpha q,
\]
\[
[-\ii\sg,W_1]_{12}=-2\ii(W_1)_{12}.
\]
所以 $(1,2)$ 元成为
\[
4\beta q_x-2\ii(W_1)_{12}+2\ii\alpha q=0.
\]
逐项移项并除以 $2\ii$, 得
\[
(W_1)_{12}=\alpha q-2\ii\beta q_x.
\]
反 Hermitian 性立即确定另一非对角元:
\[
(W_1)_{21}=-\alpha\conj q-2\ii\beta\conj{q_x}.
\]

\paragraph{$k$ 项的对角元与第一个积分常数.}
$k$ 项的矩阵方程是
\[
-(W_1)_x+[-\ii\sg,W_0]+[Q,W_1]=0.
\]
中间交换子的对角元为零, 所以
\[
\begin{aligned}
\partial_x(W_1)_{11}
&=[Q,W_1]_{11}\\
&=q(W_1)_{21}+(W_1)_{12}\conj q\\
&=q(-\alpha\conj q-2\ii\beta\conj{q_x})
+(\alpha q-2\ii\beta q_x)\conj q\\
&=-2\ii\beta(q\conj{q_x}+q_x\conj q)\\
&=-2\ii\beta\partial_x|q|^2.
\end{aligned}
\]
关于 $x$ 积分, 得
\[
(W_1)_{11}=-2\ii\beta|q|^2+\chi_1(t).
\]
要确定它, 先按原来的端态读取时间矩阵:
\[
 V\longrightarrow-F(k)U_\infty(k).
\]
该矩阵的 $(1,1)$ 元恰为 $\ii kF(k)$.
精确展开
\[
\begin{aligned}
\ii kF(k)
={}&4\ii\beta k^3-\ii\alpha k^2
-2\ii\beta q_0^2 k
+\dfrac{\ii\eta_0}{a}
+\dfrac{\ii\eta_0}{k-a}.
\end{aligned}
\]
因此 $k$ 系数在端态是 $-2\ii\beta q_0^2$.
另一方面 $q\to q_0$, 所以积分式的端态是
$-2\ii\beta q_0^2+\chi_1(t)$.
两者相等迫使 $\chi_1(t)=0$.
于是 $W_1$ 的四元已全部确定为 \eqref{eq:V1} 中的规范场表达式.

\paragraph{$k$ 项的非对角元恢复 $(W_0)_{12}$.}
同一个 $k$ 项方程的 $(1,2)$ 元为
\[
-\partial_x(W_1)_{12}-2\ii(W_0)_{12}
-2(W_1)_{11}q=0.
\]
解出未知元, 代入刚确定的 $W_1$, 得
\[
\begin{aligned}
(W_0)_{12}
&=\dfrac{\ii}{2}
\bigl(\partial_x(W_1)_{12}+2(W_1)_{11}q\bigr)\\
&=\dfrac{\ii}{2}
\bigl(\alpha q_x-2\ii\beta q_{xx}
-4\ii\beta|q|^2q\bigr)\\
&=\dfrac{\ii\alpha}{2}q_x+\beta q_{xx}+2\beta|q|^2q.
\end{aligned}
\]
反 Hermitian 性给出
\[
(W_0)_{21}
=\dfrac{\ii\alpha}{2}\conj{q_x}
-\beta\conj{q_{xx}}-2\beta|q|^2\conj q.
\]

\paragraph{常数项的对角元与第二个积分常数.}
将共振项拆成留数项与常数交换子后,
常数项为
\[
Q_t-(W_0)_x+[Q,W_0]+[\sg,R]=0.
\]
$Q_t$ 与 $[\sg,R]$ 的对角元均为零.
所以
\[
\begin{aligned}
\partial_x(W_0)_{11}
&=q(W_0)_{21}+(W_0)_{12}\conj q\\
&=q\left(\dfrac{\ii\alpha}{2}\conj{q_x}
-\beta\conj{q_{xx}}-2\beta|q|^2\conj q\right)\\
&\quad+\left(\dfrac{\ii\alpha}{2}q_x
+\beta q_{xx}+2\beta|q|^2q\right)\conj q\\
&=\dfrac{\ii\alpha}{2}(q\conj{q_x}+q_x\conj q)
+\beta(q_{xx}\conj q-q\conj{q_{xx}})\\
&=\partial_x\left[
\dfrac{\ii\alpha}{2}|q|^2
-\beta(q\conj{q_x}-q_x\conj q)
\right].
\end{aligned}
\]
在最后一步, 乘积求导的中间两项
$q_x\conj{q_x}$ 正负相消.
因此 $(W_0)_{11}$ 与方括号之间只差一个依赖 $t$ 的积分常数.
由上面的端态展开, 常数项对角元必须是
$\dfrac{\ii\eta_0}{a}$.
方括号的端态是 $\dfrac{\ii\alpha q_0^2}{2}$,
因此积分常数为
\[
\dfrac{\ii\eta_0}{a}-\dfrac{\ii\alpha q_0^2}{2}
=\dfrac{\ii}{2}\left(-\alpha q_0^2+\dfrac{2\eta_0}{a}\right)
=\dfrac{\ii\Omega}{2}.
\]
于是完整时间矩阵的两个剩余对角元为
\[
\begin{aligned}
(W_0)_{11}
&=\dfrac{\ii\alpha}{2}|q|^2
-\beta(q\conj{q_x}-q_x\conj q)+\dfrac{\ii\Omega}{2},\\
(W_0)_{22}&=-(W_0)_{11}.
\end{aligned}
\]
至此 $W_1,W_0$ 的四个独立标量项及全部积分常数均被恢复,
没有留下任意的时间函数.

\paragraph{共振留数的四个介质矩阵元.}
在有理部分利用
\[
 U(k)= U(a)-\ii(k-a)\sg,
\]
得到
\[
\dfrac{\ii([ U(k),R]-R_x)}{k-a}
=\dfrac{\ii([ U(a),R]-R_x)}{k-a}+[\sg,R].
\]
故共振留数为零即
$R_x=[-\ii a\sg+Q,R]$.
把两矩阵完整写出:
\[
-\ii a\sg+Q
=\begin{pmatrix}-\ii a&q\\-\conj q&\ii a\end{pmatrix},\qquad
R=\begin{pmatrix}\eta&-p\\-\conj p&-\eta\end{pmatrix}.
\]
两个乘积之差的独立元为
\[
\begin{aligned}
\relax[-\ii a\sg+Q,R]_{11}
&=(-\ii a\eta-q\conj p)-(-\ii a\eta+p\conj q)\\
&=-q\conj p-\conj q p,\\
[-\ii a\sg+Q,R]_{12}
&=(\ii ap-q\eta)-(\eta q-\ii ap)\\
&=2\ii ap-2\eta q.
\end{aligned}
\]
左侧对应元是 $\eta_x,-p_x$, 所以
\[
\eta_x=-(q\conj p+\conj q p),\qquad
p_x=2\eta q-2\ii ap.
\]
另两元由 Hermitian 性与迹零给出共轭方程, 不再增加独立方程.

\paragraph{常数非对角元恢复光场方程.}
$[\sg,R]_{12}=2R_{12}=-2p$,
而 $[Q,W_0]_{12}=-2(W_0)_{11}q$.
因此常数项的 $(1,2)$ 元是
\[
q_t-\partial_x(W_0)_{12}-2(W_0)_{11}q-2p=0.
\]
把已确定的分量代入:
\[
\begin{aligned}
\partial_x(W_0)_{12}+2(W_0)_{11}q
={}&\dfrac{\ii\alpha}{2}q_{xx}+\beta q_{xxx}
+2\beta(|q|^2q)_x+\ii\alpha|q|^2q\\
&-2\beta q(q\conj{q_x}-q_x\conj q)
+\ii\Omega q.
\end{aligned}
\]
非线性导数的全部乘积展开为
\[
\begin{aligned}
2\beta(|q|^2q)_x
&=2\beta(q_x\conj q+q\conj{q_x})q
+2\beta|q|^2q_x\\
&=4\beta|q|^2q_x+2\beta q^2\conj{q_x},\\
-2\beta q(q\conj{q_x}-q_x\conj q)
&=-2\beta q^2\conj{q_x}+2\beta|q|^2q_x.
\end{aligned}
\]
相加后 $q^2\conj{q_x}$ 消去, 仅剩 $6\beta|q|^2q_x$.
故
\[
q_t=\ii\alpha\left(\dfrac{q_{xx}}{2}+|q|^2q\right)
+\beta(q_{xxx}+6|q|^2q_x)+2p+\ii\Omega q.
\]
最后, 由 $q=e^{\ii\Omega t}\widetilde q$,
\[
q_t=e^{\ii\Omega t}
(\widetilde q_t+\ii\Omega\widetilde q),\qquad
q_x=e^{\ii\Omega t}\widetilde q_x,
\]
且 $p=e^{\ii\Omega t}\widetilde p$, $|q|=|\widetilde q|$.
把这些关系代入规范光场方程, 除以 $e^{\ii\Omega t}$,
两侧的 $\ii\Omega\widetilde q$ 恰好相消, 得到原 HMB 光场方程.
介质方程中空间导数不作用于共同的时间相位,
所以同样直接还原原场的两条介质方程.

\paragraph{从所得矩阵重新识别散射数据.}
在 $x\to-\infty$ 时 $M\to Y$,
所以 $Me^{\ii\theta\sg}$ 是负端归一化的 Jost 解.
在 $x\to+\infty$ 时
$M\to Y\diag(s_{22},s_{22}^{-1})$,
右乘 $\diag(s_{22}^{-1},s_{22})$ 恰好消去该端态因子.
因此正端归一化的 Jost 解是
\[
M\diag(s_{22}^{-1},s_{22})e^{\ii\theta\sg}.
\]
对角矩阵与指数对易, 故正负端解满足
\[
\phi_+=\phi_-\diag(s_{22}^{-1},s_{22}).
\]
其非对角散射元都是零.
透射系数恰为原来的两因子乘积,
而
\[
b_1=c_1s_{22}'(\zeta_1)
=\dfrac{4\ii\kappa}{1+h_0^2}
\dfrac{1+h_0^2}{4\ii\kappa}=1.
\]
由此验证重构后的实际散射数据与最初输入的数据一致.

\subsubsection*{双曲形式的指数展开与传播参数}
原记号 $s=\xi+\ii\nu$ 给出
\[
\begin{aligned}
T=e^{2s}&=e^{2\xi}e^{2\ii\nu},\\
\conj T&=e^{2\xi}e^{-2\ii\nu},\\
|T|^2&=e^{4\xi}.
\end{aligned}
\]
$P_{12}$ 的分子先按乘法展开:
\[
\begin{aligned}
(T-h_0)(1-h_0\conj T)
&=T+h_0^2\conj T-h_0(1+|T|^2).
\end{aligned}
\]
分子分母同除以严格正数 $e^{2\xi}$, 分子成为
\[
\begin{aligned}
e^{2\ii\nu}+h_0^2e^{-2\ii\nu}
-h_0(e^{2\xi}+e^{-2\xi})
={}&(1+h_0^2)\cos(2\nu)
+\ii(1-h_0^2)\sin(2\nu)\\
&-2h_0\cosh(2\xi).
\end{aligned}
\]
这里使用的三个指数恒等式分别是
\[
e^{2\ii\nu}=\cos(2\nu)+\ii\sin(2\nu),\quad
e^{-2\ii\nu}=\cos(2\nu)-\ii\sin(2\nu),
\]
\[
e^{2\xi}+e^{-2\xi}=2\cosh(2\xi).
\]
分母成为
\[
\begin{aligned}
e^{-2\xi}\Delta_-
&=(1+h_0^2)(e^{-2\xi}+e^{2\xi})
-2h_0(e^{2\ii\nu}+e^{-2\ii\nu})\\
&=2(1+h_0^2)\cosh(2\xi)-4h_0\cos(2\nu).
\end{aligned}
\]
相除即得到 \eqref{eq:hyperbolic-projector}.
因为 $\cosh(2\xi)\ge1$ 且 $\cos(2\nu)\le1$,
\[
2(1+h_0^2)\cosh(2\xi)-4h_0\cos(2\nu)
\ge2(1+h_0^2)-4h_0=2(1-h_0)^2>0.
\]
因此双曲形式与有理投影形式给出相同的全局光滑性判断.
若进一步把传播函数代入原有的 $s$, 复系数可以完全分开:
\[
\begin{aligned}
F(\ii b)
&=-4\beta b^2-\ii\alpha b-2\beta q_0^2
+\dfrac{\eta_0}{a(\ii b-a)}\\
&=-4\beta b^2-2\beta q_0^2-\dfrac{\eta_0}{a^2+b^2}
-\ii b\left(\alpha+\dfrac{\eta_0}{a(a^2+b^2)}\right).
\end{aligned}
\]
其中最后一步将共振分母乘以其共轭:
\[
\dfrac{\eta_0}{a(\ii b-a)}
=\dfrac{\eta_0(-a-\ii b)}{a(a^2+b^2)}.
\]
因此
\[
\begin{aligned}
\xi
&=\kappa\left[
x+\left(4\beta b^2+2\beta q_0^2
+\dfrac{\eta_0}{a^2+b^2}\right)t
\right],\\
\nu
&=\kappa b\left(\alpha+\dfrac{\eta_0}{a(a^2+b^2)}\right)t.
\end{aligned}
\]
这两式没有重新定义谱变量或引入简算别名;
它们只是将原先 $s$ 的实部与虚部逐项算出.
% END EXPANSION 08-J
